\documentclass[11pt]{article}
% Diplomatic transcription macros (12 Nov 2026). Typed text is plain; the handwritten layer is blue,
% deletions red, typewriter strike-outs grey, doubtful readings in brackets with a question mark.
\usepackage[utf8]{inputenc}
\usepackage[T1]{fontenc}
\usepackage{lmodern}
\usepackage[a4paper,margin=2.2cm]{geometry}
\usepackage{amsmath,amssymb}
\usepackage{xcolor}
\usepackage{tikz-cd}
\usepackage{environ}
\setlength{\parindent}{0pt}
\setlength{\parskip}{4pt}
\title{EGA V, English translation by P. Blass, 1987 (cote 24)\\[4pt]\large Diplomatic transcription from the Montpellier scans}
\author{AI-transcribed by Claude Opus 5.5}
\date{12 November 2026}
\newcommand{\scanpage}[1]{\par\bigskip\noindent{\small\sffamily[scan #1]}\par}
\newcommand{\mspage}[1]{{\small\sffamily[p.\,#1]}}
\newcommand{\hand}[1]{\textcolor{blue}{#1}}
\let\over\relax
\newcommand{\over}[2]{\textcolor{gray}{#1}\textcolor{blue}{\ensuremath{^{\!\nearrow}}#2}}
\newcommand{\inserted}[2][]{\textcolor{blue}{$\langle$#2$\rangle$}}
\newcommand{\margin}[1]{\textcolor{blue}{\{\textsf{marg.:} #1\}}}
\newcommand{\deleted}[1]{\textcolor{red}{[\textsf{del.:} #1]}}
\newcommand{\xout}[1]{\textcolor{gray}{[\textsf{x:} #1]}}
\newcommand{\uline}[1]{\underline{#1}}
\newcommand{\uuline}[1]{\underline{\underline{#1}}}
\newcommand{\unsure}[1]{[#1?]}
\newcommand{\illegible}[1]{[\textsf{illeg.\ #1}]}
\newcommand{\diagram}[1]{\par{\centering #1\par}}
\newenvironment{diagramblock}{\par\begin{center}}{\end{center}\par}

\begin{document}
\maketitle

% EGA V, English translation by P. Blass, 1987 (cote 24): archive cote 24, site page 1 (PDF page 2).
% typescript. Draft compared: none.
% AI-transcribed by Claude Opus 5.5 (claude-opus-5-5, adaptive thinking, effort high), Message Batches API, batch ega24b (msgbatch_01RGJJeHdR96U16ZWFuUjtDx), 2026-11-13; not yet checked by eye.
\scanpage{24/1}
\mspage{1}

EGA$\overline{\underline{V}}$ based on notes of
Alexandre Grothendieck

Stage 0 - translation

by

Piotr Blass

\clearpage
% EGA V, English translation by P. Blass, 1987 (cote 24): archive cote 24, site page 2 (PDF page 3).
% typescript. Draft compared: none.
% AI-transcribed by Claude Opus 5.5 (claude-opus-5-5, adaptive thinking, effort high), Message Batches API, batch ega24b (msgbatch_01RGJJeHdR96U16ZWFuUjtDx), 2026-11-13; not yet checked by eye.
\scanpage{24/2}
\mspage{2}

Summary

This formulation gives pêle mêle a detailed summary of the
set of results that should appear in a final formulation. To
arrive at the latter we need to reorganize??? thoroughly the
present stage zero the first step should probably be to make a new
plan (in which without a doubt the present sections 11, 12, 14, 15
will come much earlier). I have not even written No 16 which
should neither in principle cause any difficulty nor does it
influence in any way the previous Nos since what is involved is a
simple matter of translation.

You will notice the presence of a proposition 10.3 which
should appear in a previous paragraph.

I would like to tell you in this connection that I have
several other results quite diverse but all dealing with
birational mappings that I would love to include somewhere.

It seems to me that there is not enough to make a paragraph.
Do you have a suggestion where to put them?

I plan to send them to you soon as well as No. 16 of the
present notes. [Ask AG if No 16 has ever been written.] \hand{$^{*}$}

In addition the present section 20 will probably blow up into
two sections, one consisting of results of the type ``elementary
geometry'' on grassmonians.

If need be, could one include there also (lacking a better
place) the supplements that I told you about dealing with
birational transformations?

\margin{$^{*}$ Ask AG.}

\clearpage
% EGA V, English translation by P. Blass, 1987 (cote 24): archive cote 24, site page 3 (PDF page 4).
% typescript. Draft compared: none.
% AI-transcribed by Claude Opus 5.5 (claude-opus-5-5, adaptive thinking, effort high), Message Batches API, batch ega24b (msgbatch_01RGJJeHdR96U16ZWFuUjtDx), 2026-11-13; not yet checked by eye.
\scanpage{24/3}
\mspage{3}

\uline{Hyperplane Sections and Conic Projections}

1)\quad Notations.

2)\quad Generic hyperplane section, local properties.

3)\quad Generic hyperplane section geometric irreducibility and
connectedness.

4)\quad Variable hyperplane section: ``sufficiently general''
sections.

5)\quad Theorems of Seidenberg type.

6)\quad Connectedness of any hyperplane section.

7)\quad Application to the constructions of special hyperplane
sections and multisections.

8)\quad Dimension of the set of exceptional hyperplanes.

9)\quad Change of projective immersion.

10)\quad Pencils of hyperplane sections and fibrations of blown up
varieties.

11)\quad Grassmanians.

12)\quad Generalization of forementioned results to linear sections.

13)\quad Elementary morphisms and a theorem of M. Artin.

14)\quad Conic projections.

15)\quad Axiomatization of some of the previous results.

16)\quad Translation into the language of linear systems.

\clearpage
% EGA V, English translation by P. Blass, 1987 (cote 24): archive cote 24, site page 4 (PDF page 5).
% typescript. Draft compared: none.
% AI-transcribed by Claude Opus 5.5 (claude-opus-5-5, adaptive thinking, effort high), Message Batches API, batch pilot (msgbatch_0189Nhf1xB615YWnztubtLfc), 2026-11-12; not yet checked by eye.
\scanpage{24/4}
\mspage{4}

\uline{HYPERPLANE SECTIONS AND CONIC PROJECTIONS}

\uline{Section 1: Preliminaries and Notation.}

Let S be a prescheme, let $\underline{E}$ be a locally free module of
finite type over S let $\underline{E}^{\vee}$ be its dual. We denote by P = $\underline{P}(\underline{E})$ the
projective fibration defined by $\underline{E}$ and by P$^{\vee}$ the projective
fibration defined by E$^{\vee}$. We shall call P$^{\vee}$ the \uline{scheme of}
\uline{hyperplanes} of P. This terminology can be justified as follows.
Let $\xi$ be a section of P$^{\vee}$ over S which is therefore determined by an
invertible quotient module L of $\underline{E}^{\vee}$. From it we obtain an
invertible quotient L$_{P}$ of $(\underline{E}^{\vee})_{\underline{P}} = (\underline{E}_{P})^{\vee}$, on the other hand, we have
the invertible quotient module O$_{P}$(1) of E$_{P}$. Passing to duals we
may take $\underline{L}_{P}{-1}$ resp. $\underline{O}_{P}(-1)$ to be invertible submodules (locally
direct \uline{summands or factors}) of $\underline{E}_{P}$ (resp. of $(\underline{E}_{P})^{\vee}$) and the pairing
$EP \otimes E^{\vee}_{P} \to \underline{O}_{P}$ defines therefore a natural pairing (*) $\underline{O}_{P}(-1) \otimes \underline{L}_{P}{-1}$
$\to \underline{O}_{P}$ or also the transpose homomorphism (**) $\underline{O}_{P} \to \underline{O}_{P}(1) \otimes \underline{L}_{P} =$
$\underline{O}_{P}(1)$ i.e. a section of $\underline{L}_{P}(1)$ canonically defined by $\xi$. The
``divisor'' of that section i.e. the closed subscheme H$_{\xi}$ of $\underline{P}$ defined
by the image ideal of (*) is called the hyperplane in $\underline{P}$ defined by
the element $\xi \in$ P$^{\vee}$(S). We could describe it by noting that locally
over $\underline{S}$, $\xi$ is given by a section $\phi$ of E such that $\phi$ (s) $\neq 0$ for all
s ($\phi$ is determined by $\xi$ up to multiplication by an invertible
section of $\underline{O}_{S}$); since E = p$^{*}$(O$_{P}$(1)) (p: P $\to$ S being the
projection), $\phi$ can be considered as a section of O$_{P}$(1) the divisor
of which is nothing else but H$_{\xi}$.

Of course, if we consider L-1 as an invertible submodule of E
locally a direct factor in E then the correspondence between $\xi$
% NOTE: a white strip covers the next typed line (only faint traces show through); nothing else is typed below it on the page

\clearpage
% EGA V, English translation by P. Blass, 1987 (cote 24): archive cote 24, site page 5 (PDF page 6).
% typescript. Draft compared: none.
% AI-transcribed by Claude Opus 5.5 (claude-opus-5-5, adaptive thinking, effort high), Message Batches API, batch ega24b (msgbatch_01RGJJeHdR96U16ZWFuUjtDx), 2026-11-13; not yet checked by eye.
\scanpage{24/5}
\mspage{5}

(i.e. $\underline{L}$ or $\underline{L}$-1 \hand{\unsure{$\subset$|$\in$}} $\underline{E}$) and $\phi$ is obtained by taking for $\phi$ a
section of $\underline{L}$-1 which does not vanish at any point i.e. by a
trivialization of $\underline{L}$-1 (which exists in every case locally). Let us
note that H$\xi$ is nothing else but P($\underline{E}$ / $\underline{L}$-1) (canonical isomorphism)
that is a third way of describing H$\xi$ (N.B. P($\underline{E}$ / $\underline{L}$-1) is indeed
canonically embedded in P = P($\underline{E}$) which has the advantage of
proving in addition that H is a projective fibration over S and is
a fortiori smooth over S. (Again it would be necessary to say in
par 17 that a projective fibration is smooth). Without a doubt it
would be better to begin with that one. [Ask AG par 17 of what?]

\uline{Remarks}

The formation of H$\xi$ from $\xi$ is compatible with base change in
other words one finds a homomorphism of functors [Ask AG about
next symbol]

\quad $^{\circ}$

(Sch) /S $\rightarrow$ (Ens), PV $\rightarrow$ \uline{Div}(P/S) where the second term denotes the
functor of ``relative divisors'' of P/S whose values at S' (an
arbitrary S prescheme) is the set of closed subschemes of $\underline{P}_{S}$' which
are complete intersections transversal and of codimension 1
relative to S' (compare Par. 19).

It is easy to show that this homomorphism of functors is a
monomorphism, in other words that $\xi$ is determined of H$\xi$. (This
last fact justifies the terminology ``scheme of hyperplanes'' used
above.) We shall see that the functor \uline{Div}(P/S) is representable
by the prescheme (direct) sum of P(Symm$^{k}$(E)) so that $\underline{P}$ can be
identified to an open and closed subscheme of \uline{Div}(P/S).....

\clearpage
% EGA V, English translation by P. Blass, 1987 (cote 24): archive cote 24, site page 6 (PDF page 7).
% typescript. Draft compared: none.
% AI-transcribed by Claude Opus 5.5 (claude-opus-5-5, adaptive thinking, effort high), Message Batches API, batch ega24b (msgbatch_01RGJJeHdR96U16ZWFuUjtDx), 2026-11-13; not yet checked by eye.
\scanpage{24/6}
\mspage{6}

(N.B. to tell the truth the determination of the relative divisors
of P/S could be done with the means available right now, using
results of Ch. III and could be added as an example to Par. 19..).

Let us now make the base change $S^1 = P^V \to S$ and let us
consider the diagonal section (or ``generic section'') of $P^V{}_{S1}$ =
$\underline{P}(E^V{}_{S1})$ over $S^1$ : we find a closed subscheme $\underset{S}{\underline{H}}$ of $P_{S1} = P \times_S P^V$
which is called the incidence scheme between $\underline{P}$ and $\underline{P}^V$ defined by
the image ideal of the canonical homomorphism
$O_{\underline{P}}(-1) \otimes_S O_{\underline{P}^V}(-1) \to O_{\underline{P}\,\underset{S}{\times}\,\underline{P}^V}$

from which one sees that it is a projective fibration over $P^V$, and
by symmetry it is also a projective fibration over P.  Let us note
that one recovers the ``special'' hyperplanes $H_\xi$ (for $\xi$ a section of
$P^V$ over S) by starting out from the ``universal hyperplane'' $\underline{H}$ and by
taking its inverse image for the base change $S \xrightarrow{\xi} P^V$.

The same remark holds for every point of $P^V$ with values in an
arbitrary S-prescheme S' which (considered as a section of $\underline{P}_S{}'$ over
S') allows us to define an $H_\xi$ \hand{$\subset$} $P_S{}'$ ; the latter is nothing else
but the inverse image of H by the base change $S' \xrightarrow{\xi} \underline{P}^V$.  In what
follows we assume a prescheme X of finite type over S and an S
morphism $f{:}X \;\to P$.  One of the main objectives of this paragraph
(section) is to study for every hyperplane $H_\xi$ of P its inverse
image $\underline{Y}_\xi = f^{-1}(H_\xi) = X \underset{P}{\times} H_\xi$ and especially to relate
% NOTE: a white box covers the end of this last typed line, after "to relate"; the rest of the line is hidden.

\clearpage
% EGA V, English translation by P. Blass, 1987 (cote 24): archive cote 24, site page 7 (PDF page 8).
% typescript. Draft compared: none.
% AI-transcribed by Claude Opus 5.5 (claude-opus-5-5, adaptive thinking, effort high), Message Batches API, batch ega24b (msgbatch_01RGJJeHdR96U16ZWFuUjtDx), 2026-11-13; not yet checked by eye.
\scanpage{24/7}
\mspage{7}

the properties of X and of $Y_\xi$.  As usual one also has to consider
the $P(S^1)$, $S^1$ arbitrary S scheme, in this case $H_\xi$ is a hyperplane
in $\underline{P}_S{}^1$ and we put again
\[ Y_\xi = f_S{}'\,{}^{-1}(H_\xi) = x_S{}'\; x_{\underline{P}_{S}}{}' \qquad\qquad H_\xi = X \quad x_P \;\; H_\xi \]
where the subscript S' denotes as usual the effect of the base
change S' \qquad\qquad S and where in the last expression we consider
$H_\xi$ as a P scheme via the combined morphism $H_\xi$ \qquad $\underline{P}_S{}'$ \quad P.  It is
therefore again convenient to consider the case where $\xi$ is
``universal'' i.e. where S' = $P^V$ and $\xi$ is the diagonal section so
that H = $\underline{H}$, in this case one observes (up to better notations to
be suggested by Dieudonne) that $\underline{Y} = Y_\xi$ in the general case of a $\xi$
: S' \qquad\qquad $P^V$, one has therefore also $Y_\xi = \underline{Y}\; x_{\underline{P}^{V}} \quad S^1$.  Finally if
F is a \uline{sheaf of modules}* over $XP^V$ we denote by G its inverse image
over Y by G its inverse image over $\underline{H}$ so that one also has
\[ G_\xi = \underline{G} \quad \bigotimes_{O_{\underline{P}}{}^V} \quad O_S{}^1. \]

Let us summarize in a small diagram the essentials of the
constructions and notations considered.

\begin{diagramblock}
\begin{tikzcd}
F & & \underline{G} & G_\xi \\
X \arrow[d] & X\; x_S\; \underline{P}^V \arrow[l] \arrow[d, dash] & \underline{Y} \arrow[l] \arrow[d] & Y_\xi \arrow[l] \arrow[d] \\
\underline{P} \arrow[d] & \underline{P}\; x_S\; \underline{P}^V \arrow[l] \arrow[d, dash] & \underline{H} \arrow[l] \arrow[dl] & H_\xi \arrow[l, dash] \arrow[d] \\
S \arrow[r, dash] & \underline{P}^V & & S' \arrow[ll]
\end{tikzcd}
\end{diagramblock}

\margin{* ASK AG if \uline{Module} always means coherent or quasi-coherent sheaf of modules.}

\clearpage
% EGA V, English translation by P. Blass, 1987 (cote 24): archive cote 24, site page 8 (PDF page 9).
% typescript. Draft compared: none.
% AI-transcribed by Claude Opus 5.5 (claude-opus-5-5, adaptive thinking, effort high), Message Batches API, batch ega24b (msgbatch_01RGJJeHdR96U16ZWFuUjtDx), 2026-11-13; not yet checked by eye.
\scanpage{24/8}
\mspage{8}

(the squares and diamonds appearing in this diagram are
Cartesian). In the next section we will study systematically the
following case: $S^{1}$ is the spectrum of a field K and its image in P
is generic in the corresponding fiber $P^{V}_{\unsure{S|s}}$. After making the base
change Spec k(s) $\rightarrow$ S we are reduced to the case where S is the
spectrum of a field k, which is what we assume in the next
section. Also the majority of properties studied for X and Y are
of ``geometric nature'' and therefore invariant under base change,
which allows us also (without loss of generality) to limit
ourselves to the case where K is algebraically closed or to the
case where K = k($\eta$), $\eta$ , being the generic point of $P^{V}$ and . $\xi$ :
Spec(K) $\rightarrow$ $P^{V}$ is of course the canonical morphism. We also note
that for geometric questions concerning X, Y we can (after making
a base change) restrict ourselves to the case of k algebraically
closed. A terminological reminder. If f is an immersion we
usually call Y a hyperplane section of X (relative to the
projective immersion f and the hyperplane H). There is no reason
not to extend this terminology to the case of an arbitrary f.

\clearpage
% EGA V, English translation by P. Blass, 1987 (cote 24): archive cote 24, site page 9 (PDF page 10).
% typescript. Draft compared: none.
% AI-transcribed by Claude Opus 5.5 (claude-opus-5-5, adaptive thinking, effort high), Message Batches API, batch ega24b (msgbatch_01RGJJeHdR96U16ZWFuUjtDx), 2026-11-13; not yet checked by eye.
\scanpage{24/9}
\mspage{9}

\uline{Section 2}

\uline{20.2}

\uline{Study of a generic hyperplane section : local properties.}

Let us recall that now S=Spec(k), k is a field.  If $\eta$ is a
point of $\underline{P}V$ and if $\xi$ :Spec k($\eta$) $\rightarrow$ PV is the canonical morphism we
also \uline{put} $H_{\eta}$, $Y_{\eta}$, $G_{\eta}$ in place of $H_{\xi}$ , $Y_{\xi}$ , $G_{\xi}$ .

In this section (numero) $\eta$ denotes always the generic point
of PV.

\uline{Proposition 2.1}

Let us assume that X is irreducible.  Then $\underline{Y}_{\eta}$ is irreducible
or empty and in the first case it dominates X\uline{???} [illegible ask
AG] $\overline{Y}$ is irreducible.

Indeed since $\underline{H} \rightarrow P$ is a \uline{projective fibration} that is also
true for $\underline{Y} \rightarrow X$ which implies that $\underline{Y}$ is irreducible if X is
irreducible.  So the generic fiber Y of $\underline{Y}$ over PV is irreducible or
empty in the first case its generic point is the generic point of
$\underline{Y}$ which therefore lies over the generic point of X.  q.e.d

\uline{Proposition 2.2}

Let Z be a subset of P.  Then its inverse image $Z_{\eta}$ in $H_{\eta}$ is
empty if and only if every point of Z is closed.  In particular if
Z is constructible then $Z_{\eta} = \phi$ if and only if Z is finite.

\clearpage
% EGA V, English translation by P. Blass, 1987 (cote 24): archive cote 24, site page 10 (PDF page 11).
% typescript. Draft compared: none.
% AI-transcribed by Claude Opus 5.5 (claude-opus-5-5, adaptive thinking, effort high), Message Batches API, batch ega24b (msgbatch_01RGJJeHdR96U16ZWFuUjtDx), 2026-11-13; not yet checked by eye.
\scanpage{24/10}
\mspage{10}

\uline{Outline of proof of proposition 2.2.}*

We may suppose that Z is reduced to a single point z and we
only have to prove that the image of H in P consists exactly of
the non-closed points of P.

Denoting by X the closure of z and using 2.1 we only have to
prove that $Z_{\eta} = \phi$ if and only if X is finite (X being a closed
subscheme of $\underline{P}$).  Replacing X by $X_{k(\eta)} \to \underline{P}_{k(\eta)}$ the 'only if'
[French:  'il faut' ``one must''] part results from the following
fact for which we have to have a reference and which fact deserves
to be restated here as a lemma:  if Y is \uline{any} hyperplane section of
X and if $Y_{\eta} = \phi$ then X is finite (indeed X \quad\quad $\underline{P}$ - H is affine and
projective....).  The 'is sufficient' part one needs is immediate
for example by noting that $\underline{Y}$ is a projective fibration of relative
dimension (n - 1) over X (n being the relative dimension of $\underline{P}$ and
$\underline{P}^{V}$ over S), thus X being finite over k, $\underline{Y}$ is of absolute dimension
n - 1, , < n = dim $\underline{P}$ > thus the morphism Y $\to$ $\underline{P}^{V}$ cannot be dominant
thus its generic fiber $Y_{\eta}$ is empty.

\uline{Corollary 2.3}

Let f:X \quad\quad\quad\quad $\underline{P}$ be a morphism of finite type and let Z be
a constructible subset of X.  In order for its inverse image in $Y_{\eta}$
to be empty it is necessary and sufficient that the image f(Z)
should be finite.  In particular in order for $Y_{\eta}$ to be empty it is
necessary and sufficient that f(X) should be finite.

\margin{* Ask Grothendieck: what is the meaning and role of underlined
capital letters?}

\clearpage
% EGA V, English translation by P. Blass, 1987 (cote 24): archive cote 24, site page 11 (PDF page 12).
% typescript. Draft compared: none.
% AI-transcribed by Claude Opus 5.5 (claude-opus-5-5, adaptive thinking, effort high), Message Batches API, batch ega24b (msgbatch_01RGJJeHdR96U16ZWFuUjtDx), 2026-11-13; not yet checked by eye.
\scanpage{24/11}
\mspage{11}

\uline{Corollary 2.4}

Let $Z,Z'$ be two closed subsets of $X$ with $Z$ irreducible and let
$Z_\eta$ and $Z'_\eta$ be their inverse images in $Y_\eta$. In order to have
$Z_\eta$ \hand{$\subset$} $Z'_\eta$ it is necessary and sufficient that $f(Z)$ should be
finite or that we have $Z$ \hand{$\subset$} $Z'$.

In order that $Z_\eta = Z'_\eta$ it is necessary and sufficient that
$f(Z)$ and $f(Z1)$ should be finite or that $Z = Z'$.

This is an immediate consequence of 2.3 because we see that
$f(Z - Z\cap Z')$ can only be finite if $Z$ \hand{$\subset$} $Z'$ or if $f(Z)$ is finite
(because if we do not have $Z$ \hand{$\subset$} $Z'$ then $Z - Z\cap Z'$ is dense in $Z$
thus $f(Z - Z\cap Z')$ is dense in $f(Z)$, and if the former is finite and
thus closed, being constructible, so is also the latter.

\uline{Corollary 2.5}

To every irreducible component $X_i$ of $X$ such that $\dim \overline{f(X_i)} > 0$
we make correspond its inverse image $Y_{i\eta}$ in $Y_\eta$. Then $Y_{i\eta}$ is an
irreducible component of $Y_\eta$ and we obtain in this manner a one to
one correspondence between the set of irreducible components $X_i$ of
$X$ such that $\dim \overline{f(X_i)} > 0$ and the set of irreducible components of
$Y_\eta$.

Indeed it follows from 2.3 that $Y_\eta$ is union of $Y_{i\eta}$ defined
above and that the latter are closed and non-empty subsets of $\underline{Y}$,
they are also irreducible because of 2.1. Finally, they are
mutually without an inclusion relation because of 2.4, hence the
conclusion.

Let us notice that if $\dim X_i = d_i$ we have $\dim Y_i = d_i - 1$.

\clearpage
% EGA V, English translation by P. Blass, 1987 (cote 24): archive cote 24, site page 12 (PDF page 13).
% typescript. Draft compared: none.
% AI-transcribed by Claude Opus 5.5 (claude-opus-5-5, adaptive thinking, effort high), Message Batches API, batch ega24b (msgbatch_01RGJJeHdR96U16ZWFuUjtDx), 2026-11-13; not yet checked by eye.
\scanpage{24/12}
\mspage{12}

\uline{Proposition 2.6}

Let us suppose that for every irreducible component $X_i$ of X
we have dim $\overline{f(X_i)} > 0$ i.e. $Y_{i\eta} \neq \emptyset$ or that f is an immersion and
dim $\overline{f(X)} > 0$ [slightly illegible)(Ask AG)].

Then we have dim $\underline{Y}_\eta$ = dim X - 1.

We are reduced to the case where X is irreducible due to 2.5.
By the very construction Y is defined starting from $X_{k(\eta)}$ as the
divisor of a section of an invertible module over $X_{k(\eta)}$ (being the
inverse image of $O_{\underline{P}}(1)$).  On the other hand $X_{k(\eta)}$ is irreducible
(because X is such and $k(\eta)$ is a pure transcendental extension of
k which one should have indicated at the beginning of the N°) and
$Y_\eta \neq X_{k(\eta)}$ since the image of $Y_\eta$ in X (contrary to that of $X_{k(\eta)}$,
which is faithfully flat over X) is not equal to X, indeed it does
not contain the closed points of X because of 2.3.  It follows
that dim $Y_\eta$ = dim $X_{k(\eta)}$ - 1 = dim X - 1 (reference needed for the
last equality.) \uline{q.e.d.}

\uline{Proposition 2.7}

Let F be a quasi coherent module over X, hence $G_\eta$ over $Y_\eta$.
Let $Z_i$ be the associated prime cycles of F such that dim $\overline{f(Z_i)} > 0$.
Let $Z_{i\eta}$ be the inverse image of $Z_i$ in $Y_\eta$ then the $Z_{i\eta}$ are exactly
all the prime cycles associated to $G_\eta$ also their relation of
inclusion are the same or among the $Z_i$.

The last assertion is contained in 2.4.  On the other hand
since $Y \to X$ is a projective fibration, thus flat with fibers $(S_1)$
and irreducible it follows from par. 3* that the associated prime

\margin{* Ask Grothendieck about this reference.}

\clearpage
% EGA V, English translation by P. Blass, 1987 (cote 24): archive cote 24, site page 13 (PDF page 14).
% typescript. Draft compared: none.
% AI-transcribed by Claude Opus 5.5 (claude-opus-5-5, adaptive thinking, effort high), Message Batches API, batch ega24b (msgbatch_01RGJJeHdR96U16ZWFuUjtDx), 2026-11-13; not yet checked by eye.
\scanpage{24/13}
\mspage{13}

cycles to the inverse image G of F over $\overline{Y}$ are the inverse images
of the associated prime cycles of F.  Hence they are induced on
the generic fiber $Y_{\eta}$ of $\underline{Y}$ over $\check{\underline{P}}$, the fact that the associated
prime cycles to $G_{\eta}$ are the (2) \uline{non-empty} inverse images of the $Z_{i}$
which proves 2.7 by means of 2.3.

To tell the truth the passage through $\underline{Y}$ is unnecessary
(useless), we can use directly the fact that $Y_{\eta} \to X$ is flat with
fibers $(S_{1})$ (and also geometrically regular i.e. the morphism is
regular) and with irreducible fibers (and even geometrically
irreducible: it is the localization of projective schemes) remark
analogous to \uline{illegible} of 2.1.  [illegible ask AG.]

\uline{Proposition 2.8}

Let F be coherent over X and let $y \in \underline{Y}$, x is its image in X.
Let $\underline{P}$(M) be one of the following properties for a module of finite
type M over a local noetherian ring A:

(i) \quad coprof $M \leq N$ \qquad (ref)

(ii) \quad M satisfies $(S_{k})$ \qquad (ref)

(iii) \quad M is Cohen-Macaulay

(iv) \quad M is reduced \qquad (ref)

(v) \quad M is integral \qquad (ref)

Then for $H_{\eta}$ to satisfy the property $\underline{P}$ it is necessary and
sufficient that $F_{x}$ should satisfy it.

This follows immediately from results of paragraph 6* taking
into account that $Y_{\eta} \to X$ is a regular morphism so that
$O_{X,x} \to O_{Y_{\eta},y}$ should be regular.  Taking into account 2.3 we obtain
thus:

\clearpage
% EGA V, English translation by P. Blass, 1987 (cote 24): archive cote 24, site page 14 (PDF page 15).
% typescript. Draft compared: none.
% AI-transcribed by Claude Opus 5.5 (claude-opus-5-5, adaptive thinking, effort high), Message Batches API, batch ega24b (msgbatch_01RGJJeHdR96U16ZWFuUjtDx), 2026-11-13; not yet checked by eye.
\scanpage{24/14}
\mspage{14}

\uline{Corollary 2.9}

With the notations of 2.8 let $Z$ be the set of $x \in X$ such that
$\underline{P}(F_x)$ is not satisfied. Then in order for $G_\eta$ to satisfy the
condition $\underline{P}$ at all the points it is necessary and sufficient that
$f(Z)$ should be a finite subset of $\underline{P}$, or that $\overline{\mathrm{dim}}\, f(Z) = 0$.

Indeed 2.8 tells us that $h(Z)$ is a $\underline{P}$-singular subset of $G_\eta$ and
it is empty if and only if $f(Z)$ is finite by 2.3 (N.B. $h$ denotes
the morphism $Y_\eta \to X$;I have just realized that the letter $\underline{P}$ in 2.8
has been used doubly).

\uline{Corollary 2.10}

Condition for $Y_\eta$ to be reduced respectively locally integral.

\uline{Corollary 2.11}

Let $y \in \underline{Y}_\eta$, in order that $Y_\eta$ should be regular respectively
should satisfy the property $R_k$ (reference) at $y$ (respectively
should be normal at $y$) it is necessary and sufficient that $X$
should satisfy the same property at $x$. Let $Z$ be the set of those
points of $X$ where $X$ is not regular resp. $R_k$ (resp. normal); for $Y_\eta$
to be regular resp. $R_k$ (resp. normal) it is necessary and
sufficient that $f(Z)$ should be finite i.e. $\dim \overline{f(Z)} = 0$.

The same proof as 2.8 and 2.9. One must give the different
references assuring that $Z$ is closed (because we already know that
it is constructible by applying 2.3).

\clearpage
% EGA V, English translation by P. Blass, 1987 (cote 24): archive cote 24, site page 15 (PDF page 16).
% typescript. Draft compared: none.
% AI-transcribed by Claude Opus 5.5 (claude-opus-5-5, adaptive thinking, effort high), Message Batches API, batch ega24b (msgbatch_01RGJJeHdR96U16ZWFuUjtDx), 2026-11-13; not yet checked by eye.
\scanpage{24/15}
\mspage{15}

Let us note that in 2.10 and 2.11 we do not talk at all about
the corresponding geometric properties; the results anticipated
are of 'absolute' nature.  We now examine the properties of
geometric nature. (One could if one wanted to to take the
opportunity to change the n$^\circ$)

\uline{Theorem 2.12}

Suppose that $f{:}X \to \underline{P}$ is unramified.  Let $y \in Y_{\eta}$, $x$ its image
in X.  In order for X to be smooth over k at x it is ncessary and
sufficient that $Y_{\eta}$ should be smooth over $k(\eta)$ at y.

We may assume that k is algebraically closed.  If $Y_{\eta}$ is smooth
over $k(\eta)$ at y it is regular, thus since $Y_{\eta}$ is flat over X, X is
regular at x (reference), therefore it is smooth over k at x since
k is algebraically closed and thus perfect (reference).

For the converse we can (after replacing X by an open
neighborhood of x) assume that X is smooth, and (due to the
jacobian criterion of smoothness) to be defined in an open subset
U of $\underline{P}$ by $\underline{p}$ equations as $X = V(f_1,\ldots,f_{\underline{p}})$, where the differentials
$df_1,\ldots,df_{\underline{p}}$ are everywhere linearly independent.  By introducing
the affine coordinates $S_1,\ldots,S_n$ in $\underline{P}$v and the affine coordinates
$T_1$, $T_2,\ldots,T_n$ in a neighborhood of x (by choosing a hyperplane $H\infty$
[at infinity] not containing x) $y \to U_{k(\eta)}$ is then equal to
$V(f_1,\ldots,f_{\underline{p}}, \Sigma\, S_iT_i - 1)$ and it suffices to verify that the
differentials (relative to $k(\eta)$) of $f_1,\ldots,f_{\underline{p}}$, $(\Sigma\, S_iT_i) - 1)$ are
linearly independent.  However these differentials are nothing
else but the sections of $(\Omega 1_{u/k}) \otimes_k k(\quad)$ [illegible]* as follows:
$df_1,\ldots,df_{\underline{p}}$, $\Sigma\, S_i dT_i$.  Since the $df_i$ are linearly independent of

\margin{* Ask A.G.}

\clearpage
% EGA V, English translation by P. Blass, 1987 (cote 24): archive cote 24, site page 16 (PDF page 17).
% typescript. Draft compared: none.
% AI-transcribed by Claude Opus 5.5 (claude-opus-5-5, adaptive thinking, effort high), Message Batches API, batch ega24b (msgbatch_01RGJJeHdR96U16ZWFuUjtDx), 2026-11-13; not yet checked by eye.
\scanpage{24/16}
\mspage{16}

every point of U and since the $dT_i$ form a basis of $\Omega 1_{u/s}$ at every
point of U and a fortiori a system of generators we conclude
immediately the linear independence of the written down quantities
at every point of $U_{k(\eta)}$ at least when $p \leq n$, i.e.
if $E = \Omega 1_{u/k} \;/ \sum\limits_{1 \leq i \leq n} O_u df_i \neq 0$ this is a small lemma about a family of
generators of a non-zero locally free module E, thus $\sum S_i a_i$
considered as a section of $E \otimes_k k(\eta)$ does not vanish at any point.
On the other hand the case $p = n$ is trivial because then $Y_\eta = \emptyset$
% NOTE: a white rectangle covers the rest of this line after "Y_eta = (empty set)", and another white rectangle covers the beginning of the next line; any text under them is hidden.

\uline{Corollary 2.13}

Let Z be the set of points of X where X is not smooth over k.
In order that Y should be smooth over $K(\eta)$ it is necessary and
sufficient that Z should be finite. In particular if X is smooth
the same is true about Y.

Follows from 2.12 and 2.3. More generally we obtain:

\uline{Theorem 2.14}

Let y be a point of $Y_\eta$,x its image in X. Let $\underline{P}(A,K)$ be one of
the following properties for an algebra A local and noetherian
over a field K:

(i) A is geometrically regular over K.

(ii) A is geometrically $(R_k)$ over K.

(iii) A is separable over K

(iv) A is geometrically normal over K

\clearpage
% EGA V, English translation by P. Blass, 1987 (cote 24): archive cote 24, site page 17 (PDF page 18).
% typescript. Draft compared: none.
% AI-transcribed by Claude Opus 5.5 (claude-opus-5-5, adaptive thinking, effort high), Message Batches API, batch ega24b (msgbatch_01RGJJeHdR96U16ZWFuUjtDx), 2026-11-13; not yet checked by eye.
\scanpage{24/17}
\mspage{17}

then $\underline{P}(O_{X,x}, k) \Leftrightarrow \underline{P}(O_{Y_\eta,y}, k(y))$. Indeed taking into account par.
6* (iii) and (iv) follow from (ii) and 2.8 (ii). On the other
hand (i) has been proven in 2.12 and it remains to deduce (ii)
from (i). To do this let Z be the set of points where X is not
smooth over k, its inverse image $Z_\eta$ in $Y_\eta$ is therefore (by 2.12)
the set of points of $Y_\eta$ at which $Y_\eta$ is not smooth over $K(\eta)$. But
the codimension of Z in X at x is equal to that of $Z_\eta$ in $Y_\eta$ at y
because of flatness (reference par. 6*). Therefore one is $>k$ if
the other one is such which completes the proof.

\uline{Corollary 2.15}

Let Z be the set of points of X at which X is not smooth over
k (respectively is not geometrically $R_k$ over k, respectively is not
separable over k, respectively is not normal over k). In order
for $Y_\eta$ to be smooth (respectively geometrically $R_k$, respectively
separable, respectively geometrically normal) over $k(\eta)$ it is
necessary and sufficient that Z should be finite.

From writing up point of view statements 2.14 and 2.15
contain 2.12 and 2.13 (which we could thus dispense with stating
separately) on the other hand the corollary is practically more
important than the theorem which one could give in a proposition
or a preliminary lemma so that the corollary would be more
glorified. We can modify the case of property (iii):

\uline{Corollary 2.16}

(Let us still suppose that f is an immersion and also that F
is coherent), with the conditions of 2.7, in order that $Z_i$ should

\margin{* Ask Grothendieck}

\clearpage
% EGA V, English translation by P. Blass, 1987 (cote 24): archive cote 24, site page 18 (PDF page 19).
% typescript. Draft compared: none.
% AI-transcribed by Claude Opus 5.5 (claude-opus-5-5, adaptive thinking, effort high), Message Batches API, batch pilot (msgbatch_0189Nhf1xB615YWnztubtLfc), 2026-11-12; not yet checked by eye.
\scanpage{24/18}
\mspage{18}

not be immersed it is necessary and sufficient that $Z_i$ should be
such [PB incomprehensible].* If that is so then the radicial
multiplicity of $F$ at $Z_i$ at $k$ is equal to that of $G_\eta$ at $Z_i$ relative
to $k(\eta)$.

The first assertion is contained in the last assertion of
2.7. For the second, since $Y_\eta \rightarrow X$ is flat therefore the functor $F$
$\rightarrow G_\eta$ is exact, we are reduced by restriction to a neighborhood of
the general point of $Z_i$ and by a devissage (unscrewing) to the case
where $F = O_{Z_i}$ and we may assume $Z_i = X$. Also we could start by
assuming that $X$ is separable over $k$ is reduced to the case of $k$
algebraically closed [PB: incomprehensible].* Then the assertion
is contained in 2.15. (iii). Then we conclude as usual:

\uline{Corollary 2.17}

Let $Z$ be the set of points of $X$ where $F$ is not separable over
$k$ (reference). Then $G_\eta$ is separable over $k(\eta)$ if and only if $Z$ is
finite. In particular if $F$ is separable over $k$ then $G_\eta$ is
separable over $k(\eta)$.

\uline{Remark 2.18}

The key result 2.12 (and its consequences 2.13 and 2.17)
become false if we abandon the assumption that $f$ is an immersion
as we see for example by taking for $X$ a purely inseparable
covering of $\underline{P}$. However if $k$ is of characteristic zero 2.12 and
2.17 are valid without assuming that $f$ is an immersion.

\clearpage
% EGA V, English translation by P. Blass, 1987 (cote 24): archive cote 24, site page 19 (PDF page 20).
% typescript. Draft compared: none.
% AI-transcribed by Claude Opus 5.5 (claude-opus-5-5, adaptive thinking, effort high), Message Batches API, batch ega24b (msgbatch_01RGJJeHdR96U16ZWFuUjtDx), 2026-11-13; not yet checked by eye.
\scanpage{24/19}
\mspage{19}

Indeed it suffices to verify this for 2.12 and this follows
from 2.11 and from the fact that for an algebraic prescheme in
characteristic zero $\rightarrow$ smooth $=$ regular.

\uline{Section 3} (or 20.3)*
\uline{Generic hyperplane section} : \uline{geometric irreducibility and}
\uline{connectedness.}
\uline{Theorem 3.1} (Bertini - Zariski)

Let us assume dim $f(X) \geq 2$ and that X is geometrically
irreducible. Then the generic hyperplane section $\underline{Y}_{\eta}$ has the same
property.

Let K/k be the function field of X and let n = dim $\underline{P}$;
introducing the affine coordinates $T_1,\dots,T_n$ in $\underline{P}$ (by choosing a
hyperplane at infinity H$\infty$ such that f(X) is not contained in it)
and $S_1,\dots,S_n$ the affine coordinates in $\underline{P}$v, we see that the
function field L of $Y_{\eta}$ can be identified with the field of
fractions of the integral domain $K[S_1,\dots,S_n] \,/\, (\Sigma\, t_iS_i - 1)$ where
the $t_i \in K$ are the images of $T_i$ with the aid of $f{:}X \rightarrow \underline{P}$. Since
dim $\overline{f(X)} \geq 0$, the $t_i$ are not all algebraic over k, á fortiori they
are not all zero let us assume for example that $t_n \neq 0$. We realize
then immediately that we have $L = K(S_1,\dots,S_n)$ (pure transcendental
extension), $S_n \in L$ being given by the equation $\Sigma\, t_iS_i - 1 = 0$
[en fonction Fr] as a function of the $S_i$ $(1 \leq i \leq n - 1)$ and the $t_i$
$(1 \leq\ i \leq n)$. On the other hand $k^1 = k(\eta)$ can be identified with
$k(S_1,\dots,S_n)$ and the canonical inclusion $k^1 \rightarrow L$ can be obtained by
sending $S_i$ to $S_i$ [PB: check this!] i.e. $k^1$ as a subextension of L

\margin{* Ask Grothendieck about the numbering system.}

\clearpage
% EGA V, English translation by P. Blass, 1987 (cote 24): archive cote 24, site page 20 (PDF page 21).
% typescript. Draft compared: none.
% AI-transcribed by Claude Opus 5.5 (claude-opus-5-5, adaptive thinking, effort high), Message Batches API, batch ega24b (msgbatch_01RGJJeHdR96U16ZWFuUjtDx), 2026-11-13; not yet checked by eye.
\scanpage{24/20}
\mspage{20}

is the subextension generated by the $S_i$ ($1 \leq i \leq n$) or what is
evidently the same by the $S_i$ ($1 \leq i \leq n - 1$) and by $S_n = a_o + a_i\ S_1$
$+\ldots+a_{n-1}\ S_{n-1}$, $a_o = t_n{}^{-1}$, $a_i = t_it_n{}^{-1}$ for $1 \leq i \leq n - 1$.

Let us note that the field generated by the $a_i$ and by the $t_i$
is obviously the same, their common transcendence degree is
nothing else but the dimension of $f(X)$.

(N.B. It would be appropriate to include this birational
description at least as a corollary to 2.1).  The proof of 3.1 is
thus reduced to that of

\uline{Lemma 3.1.1 (ZARISKI)}

Let k be a field, K an extension of finite type over k, m an
integer $\geq 0$ $a_i$ ($0 \leq i \leq m$) the elements of K such that the
transcendence degree of $k(a_o,\ldots,a_m)$ over k is $\geq 2$, and $k^1$ the
subfield $k^1 = k(S_1,\ldots,S_m,\ a_o + \sum_{I}^{M} a_iS_i)$ of L (the $S_i$ being
indeterminates).  If K is a primary extension of k then L is a
primary extension of $k^1$.   .

This lemma or lemmata that resemble it like a brother, wander
almost everywhere in the literature, that is why I leave it up to
you:  the choice of the place from where you will copy a proof I
do not feel inspired to find a proof with my own means.

\uline{Corollary 3.2}

Assume that f is unramified or that the characteristic of k
is zero, and that the dim $\overline{f(X)} \geq 2$. Then if X is geometrically

\clearpage
% EGA V, English translation by P. Blass, 1987 (cote 24): archive cote 24, site page 21 (PDF page 22).
% typescript. Draft compared: none.
% AI-transcribed by Claude Opus 5.5 (claude-opus-5-5, adaptive thinking, effort high), Message Batches API, batch ega24b (msgbatch_01RGJJeHdR96U16ZWFuUjtDx), 2026-11-13; not yet checked by eye.
\scanpage{24/21}
\mspage{21}

integral the same is true about $Y_\eta$.  Indeed geometrically integral
$=$ geometrically irreducible $+$ separable.

\uline{Corollary 3.3}

Let us assume that $k$ is algebraically closed and that for
every irreducible component $X_i$ of $X$ we have $\dim f(X_i) \geq 2$, and
suppose also that $X$ is $\sigma$-connected where $\sigma$ is the set of closed
subsets $Z$ of $X$ such that $\dim \overline{f(X)} = 0$ (i.e. for every such $Z$,
$X - Z$ is connected).  Under such conditions $Y_\eta$ is geometrically
connected over $K(\eta)$.

Indeed by a lemma that ought to appear in par. 6* with
Hartshorne's theorem, the hypothesis signifies that we can join
any two irreducible components $X'$ and $X''$ of $X$ by a chain of
irreducible components $X_0$ and $X',\ldots,X_n = X''$ such that two
consecutive ones have an intersection $\notin \sigma$ so that the inverse
images $X'_\eta$ and $X''_\eta$ are joined by a chain of $X_{i\eta}$ which are
geometrically connected over $k(\eta)$ by 3.1 and the intersection of
two consecutive ones being $\neq \emptyset$ by 2.3.

It follows (since $Y_\eta = X_\eta$ is the union of the $X_{i\eta}$, $X_i$ running
through the set of irreducible components of $X$) that $Y_\eta$ is
geometrically connected over $k(\eta)$ q.e.d.

\uline{Section 4}\quad \uline{Study of a Variable Hyperplane Section:  ``Sufficiently}
\uline{General'' Sections}.

We return to the general situation of Section 1, S being an
arbitrary prescheme.  Also, we suppose that $X$ is of finite
presentation over $S$.

\margin{* ASK A.G.}

\clearpage
% EGA V, English translation by P. Blass, 1987 (cote 24): archive cote 24, site page 22 (PDF page 23).
% typescript. Draft compared: none.
% AI-transcribed by Claude Opus 5.5 (claude-opus-5-5, adaptive thinking, effort high), Message Batches API, batch ega24b (msgbatch_01RGJJeHdR96U16ZWFuUjtDx), 2026-11-13; not yet checked by eye.
\scanpage{24/22}
\mspage{22}

In general let us note that if $\underline{P}(Z,k)$ is a ``constructible''
property of an algebraic \uline{prescheme} Z over a field k then the set
of $\xi \in \underline{P}$v such that we have $\underline{P}(Y_\xi, k(\xi))$ is \uline{constructible} as we see
by noting that $Y_\xi$ is the fiber over $\xi$ of $\underline{Y} \to \underline{P}$v which is a morphism
of finite presentation and applying \uline{Par. 9}.* We have an analogous
remark for a property $\underline{P}(Z, M, k)$ where Z and k are as above and M
is a coherent module over Z: if G is in addition of
finite presentation over X then the set of $\xi \in \underline{P}$v such that we have
$\underline{P}(Y_\xi, G_\xi, k(\xi))$ is constructible. On the other hand in the
following sections we have developed in various cases a criterion
for the preceding set E to contain the generic point of $\underline{P}$v, S being
the spectrum of field k; \uline{being constructible} it follows that E
contains a non-empty open set this is the passage of a conclusion
from generic hyperplane section to the analogue conclusion for
``sufficiently general'' hyperplane sections.

Let us note in addition that in the case S = Spec(k) we also
have a converse: in order that the generic hyperplane section
should have the property $\underline{P}$ it is necessary and sufficient that the
$Y_\xi$ for $\xi$ in a suitable neighborhood of $\eta$ should satisfy it and it
suffices to require for $\xi$ closed in $\underline{P}$v (which for $\xi$ k algebraically
closed \uline{leads} or \uline{reduces} to considering k-rational points i.e.
hyperplane sections of X itself (without an extension prior to the
base field.)

This follows indeed from the constructibility of E and from
the fact that $\underline{P}$v is a Jacobson scheme.

Let us give as an example some special cases where the
previous considerations apply:

\margin{* \unsure{Ask} A.G.}

\clearpage
% EGA V, English translation by P. Blass, 1987 (cote 24): archive cote 24, site page 23 (PDF page 24).
% typescript. Draft compared: none.
% AI-transcribed by Claude Opus 5.5 (claude-opus-5-5, adaptive thinking, effort high), Message Batches API, batch ega24b (msgbatch_01RGJJeHdR96U16ZWFuUjtDx), 2026-11-13; not yet checked by eye.
\scanpage{24/23}
\mspage{23}

\uline{Proposition 4.2}

Let G be a module of finite presentation over X.

Let $\underline{P}$ be one of the following properties for a module M over
an algebraic scheme Z over a field K;

(i) coprof (M) $\leq$ n.

(ii) M satisfies ($S_k$).

(iii) M is Cohen Macauley.

(iv) M is without embedded prime cycle components.

(v) M is separable over K.

With these notations if E denotes the set of $\xi \in \underline{P}$ such that $G_\xi$
satisfies $\underline{P}$ then we have: (a) E is a constructible subset of $\underline{P}$v.
(b) Let us suppose that S = Spec(k), k a field, and that F is
satisfied at $\underline{P}$. Let us also suppose that in the case (v) that k
is of characteristic \unsure{O|0} \uline{or} that $f\colon X \to \underline{P}$ is unramified, then E
contains an open and dense subset of $\underline{P}$v.

\uline{Proof}

(a) follows from the fact that P is a constructible property
which we have seen in Par. 9. As to (b) it follows from the
correspsonding results of the previous two sections.

\uline{Regret}

To (b): suppose more generally that if Z is the set of
points of X where F does not satisfy $\underline{P}$, we have f(Z) is finite
i.e. dim $\overline{f(Z)} \leq 0$.

\clearpage
% EGA V, English translation by P. Blass, 1987 (cote 24): archive cote 24, site page 24 (PDF page 25).
% typescript. Draft compared: none.
% AI-transcribed by Claude Opus 5.5 (claude-opus-5-5, adaptive thinking, effort high), Message Batches API, batch ega24b (msgbatch_01RGJJeHdR96U16ZWFuUjtDx), 2026-11-13; not yet checked by eye.
\scanpage{24/24}
\mspage{24}

\uline{Proposition 4.3}

Let $\underline{P}$ be one of the following properties (for an algebraic
prescheme over a field K):

(i) Z is smooth over K.

(ii) Z satisfies the geometric property (R$_{k}$) over K.

(iii) Z is separable over K.

(iv) Z is geometrically normal over K.

(v) Z is geometrically integral over K.

(vi) Z is geometrically irreducible over K.

Let E be the set of $\xi \in \underline{P}$v such that Y$_{\xi}$ satisfies $\underline{P}$. Then (a) E is
a constructible subset of $\underline{P}$v. (b) Let us suppose S = Spec k, k a
field and let us suppose in the cases (i) to (v) that k is of
characteristic zero and that f:X $\to$ $\underline{P}$ is unramified.
Finally suppose in the cases (v) and (vi) that dim $\overline{\text{f(X)}}$ $\geq$ 2.
Assume that X satisfies $\underline{P}$ then E contains a dense open subset of
$\underline{P}$v.

\uline{Proof}

Proof is identical to that of 4.2. Let us remark that in the
cases (i) to (v) it suffices to assume \uline{only} (in b)) that f(Z) is
finite where Z is the set of points of X where $\underline{P}$ fails.

\uline{Corollary 4.4}

Let us consider the property (C$_{m}$) ``$\overline{\text{Z}}$ cannot be disconnected
by a closed subset of dimension $\leq$ m (where $\overline{\text{Z}}$ is Z x$_{K}$ $\overline{\text{K}}$, $\overline{\text{K}}$ the

\clearpage
% EGA V, English translation by P. Blass, 1987 (cote 24): archive cote 24, site page 25 (PDF page 26).
% typescript. Draft compared: none.
% AI-transcribed by Claude Opus 5.5 (claude-opus-5-5, adaptive thinking, effort high), Message Batches API, batch ega24b (msgbatch_01RGJJeHdR96U16ZWFuUjtDx), 2026-11-13; not yet checked by eye.
\scanpage{24/25}
\mspage{25}

algebraic closure of K).

Let E be the set of $\xi \in \underline{P}$v such that Y$_{\xi}$ over K($\xi$) satisfies
C$_{m}$. Then: (a) E is constructible. (b) Let us suppose that S =
Spec k, k a field, and that for every irreducible component X$_{i}$ of X
we have dim $\overline{f(X)} \geq 2$. Then if X over k satisfies C$_{m+1}$ then E$_{\xi}$
contains a dense open subset U of $\underline{P}$.

The constructibility is done by A Q T ?????? (this is a fact
that one has forgotten in Par. 9,* that perhaps it would still be
possible to repair x (or fix); the part (b) follows in principle
from 3.3 except that 3.3 has been announced in an exceedingly
special manner and consequently should be generalized (the proof
given being otherwise essentially unchanged). Also there is an
error in the statement of 4.4. Which is not valid as such if f is
quasi-finite; in the general case instead of considering the
dimension of the closed subsets of X respectively of Y$_{\xi}$ it is
sufficient to consider the dimension of their images in $\underline{P}$.
respectively in H$_{\xi}$:
Redactor demerdetur. [Latin]

\uline{Section 5} \uline{Theorems of Seidenberg Type}
\uline{5.1}

In the present section (or N$^{\circ}$) we give conditions under which
the set E defined in 4.1 is open. We deal here with properties $\underline{P}$
of local nature over X respectively Y$_{\xi}$ such that we can define the
set U of $y \in \underline{Y}$ so that (if $\xi$ denotes the image (PB ???)** of y in

\margin{* What is AQT ? Where is par 9? Ask Grothendieck (referred to
in Corollary 4.4)}
\margin{** illegible - ask AG.}

\clearpage
% EGA V, English translation by P. Blass, 1987 (cote 24): archive cote 24, site page 26 (PDF page 27).
% typescript. Draft compared: none.
% AI-transcribed by Claude Opus 5.5 (claude-opus-5-5, adaptive thinking, effort high), Message Batches API, batch ega24b (msgbatch_01RGJJeHdR96U16ZWFuUjtDx), 2026-11-13; not yet checked by eye.
\scanpage{24/26}
\mspage{26}

$\underline{P}v$) $Y_{\xi}$ satisfies $\underline{P}$ at the point y (respectively $G_{\xi}$ satisfies
condition $\underline{P}$ at y).  We give first of all the criteria for U to be
open by using paragraph 12.  As always we have $E = \underline{P} - h(\underline{Y} - U)$ it
follows that if U is open and X is proper over S (since h is
proper and a fortiori closed) then E is also open.

\uline{5.2}

As we have seen in No. 1 $\underline{Y}$ is defined in $X\ x_S\ Pv = X_{Pv}$ as the
``divisor'' of a section {\o} of $O_X(1) \otimes_S O_P(1)$ which induces for every
$\xi \in \underline{P}v$ a section ${\o}_{\xi}$ of $O_X(1) \otimes_{k(S)} O_P(1)(\xi)$ (a sheaf by the way
isomorphic non-canonically to $O_X(1) \otimes_{k(S)} k(\xi) = O_{X_{k(s)}}(1)$ )
such that $Y_{\xi}$ is nothing else but the ``divisor'' of this section
(N.B. the divisor of a section {\o} of an invertible module L is
defined as the closed subscheme defined by the image ideal of
${\o}^{-1}$: $\underline{L}^{-1} \to O$.  If $\underline{F}$ is a sheaf of modules over X then its inverse
image over $\underline{Y}$ i.e. the inverse image of $F \otimes_O O_{Pv} = F_{Pv_{S}}$ over
the subscheme Y of $X_{Pv}$ is nothing else but the cokernel of the
homomorphism ${\o}^{-1} \otimes \mathrm{id}_{\underline{F}_{P^v}}$ : $F_{Pv}(-1, -1) \to F_{Pv}$ where the
notation (-1, -1) explains itself as Mike* says.  Also $G_{\xi}$ is the
cokernel of analogous homomorphism $F_{k(\xi)}\ (-1, -1) \to F_{k(\xi)}$ where $\xi$ is
a point of P (and also we have a corresponding interpretation if $\xi$,
instead of being a point of Pv, denotes a point of Pv with values
in an S' over S).?)

In general if L is an invertible module somewhere, {\o} a
section defining the subprescheme V({\o}) then for every module F the

\clearpage
% EGA V, English translation by P. Blass, 1987 (cote 24): archive cote 24, site page 27 (PDF page 28).
% typescript. Draft compared: none.
% AI-transcribed by Claude Opus 5.5 (claude-opus-5-5, adaptive thinking, effort high), Message Batches API, batch ega24b (msgbatch_01RGJJeHdR96U16ZWFuUjtDx), 2026-11-13; not yet checked by eye.
\scanpage{24/27}
\mspage{27}

inverse image of F \uline{in} V($\varnothing$) can be identified, by the usual abuse
of language to the cokernel of id$_{\mathrm{F}}\otimes$ ? *: F $\otimes$ L-1 $\rightarrow$ F.

We say that $\varnothing$ is F regular if the preceding homomorphism is
injective.  If we choose an isomorphism of F and O$_{\mathrm{X}}$, which is
possible locally such that $\varnothing$ is identified to a section of O$_{\mathrm{X}}$ this
terminology is compatible with the one that was already introduced
elsewhere.

\uline{Proposition 5.3}

With the previous notations let U be the set of x $\in$ X$_{\mathrm{P}}$ with
image $\xi$ in Pv such that $\varnothing_{\xi}$ should be F$_{\mathrm{K}(\xi)}$ regular at x.  Then (a)
If F is of finite presentation and flat relative to S then U is
open and $\underline{\mathrm{G}}$/U is flat relative to Pv.  (b) For every s $\in$ Pv if $\eta$
denotes the generic point of P$_{\mathrm{s}}$v then U contains X$_{\mathrm{k}(\eta)}$.

\uline{Proof}

(a) Since F$_{\mathrm{P}}$v is of finite presentation and flat relative to
Pv the conclusion follows from 11.3 (and also from O$_{\mathrm{III}}$... in the
case of locally noetherian S).

(b) We may suppose S = Spec k.  The associated cycles \uline{to} F$_{\mathrm{k}(\eta)}$
are (because of Par. 3)* the inverse images of associated cycles
Z$_{\mathrm{i}}$ to F.  If f(Z$_{\mathrm{i}}$) is finite then by 2.3 Z$_{\mathrm{ik}(\eta)}$ $\cap$ Y = $\varnothing$ in the
contrary case by 2.6 for example we also have Z$_{\mathrm{ik}(\eta)}$ $\cap$ $\underline{\mathrm{Y}}$ = Z$_{\mathrm{ik}(\eta)}$
(by reason of dimension; 2.3 which was already involved in 2.6 and
is without a doubt a better reason) which proves that $\varnothing$ does not
vanish over any of the Z$_{\mathrm{ik}(\eta)}$ and therefore proves (b).

\margin{* illegible, ask A.G. or figure out.}

\clearpage
% EGA V, English translation by P. Blass, 1987 (cote 24): archive cote 24, site page 28 (PDF page 29).
% typescript. Draft compared: none.
% AI-transcribed by Claude Opus 5.5 (claude-opus-5-5, adaptive thinking, effort high), Message Batches API, batch ega24b (msgbatch_01RGJJeHdR96U16ZWFuUjtDx), 2026-11-13; not yet checked by eye.
\scanpage{24/28}
\mspage{28}

\uline{Corollary 5.4}

Let V be the set of $\xi \in P^{\vee}$ such that $\varnothing_{\xi}$ is $F_{k(\xi)}$ regular. If F
is of finite presentation then V is constructible and it contains
the generic points of the fibers of $P^{\vee}$ over S on the other hand if
also X is proper over S and P is flat over S then the set V is
open.

\uline{Remark 5.5}

Let $\xi \in P^{\vee}$ over $s \in S$ and let us suppose that $F_{s}$ should be
without associated embedded cycles. Then we see immediately that $\xi$
$\in$ V (notation of 5.4) which means also that every irreducible
component of supp $F_{k(\xi)}$ does not lie over $H_{\xi}$ (and a little less
evidently in this criterion we replace $k(\xi)$ by an arbitrary
extension of $k(\xi)$.

Let us note that the hypothesis $(S_{1})$ about $F_{s}$ which we have
just made is satisfied notably if we suppose $F_{s}$ Cohen-Macauley (a
fortiori if F is CM over S); also in this case $G_{s}$ is CM (since
locally it is deduced from $F_{k(s)}$ which is such by dividing by a\hand{$\Phi$}$\cdot$
$F_{k(s)}$ where $\varnothing$ is $F_{k(s)}$ regular). The same remarks anyway should and
will have to all be made locally above to characterize the points
of U (in place of those of V).

Using now 12.1.1* and 12.1.4* we obtain (\uline{locate reference})

\uline{Theorem 5.6}

Let us assume that F is of finite presentation flat relative
to S. Let P be one of the properties (i) to (viii) of 12.1.1 or
(if we assume F = $O_{X}$) one of the properties (i) to (iv) of 12.1.4.

\margin{* Ask AG about reference.}

\clearpage
% EGA V, English translation by P. Blass, 1987 (cote 24): archive cote 24, site page 29 (PDF page 30).
% typescript. Draft compared: none.
% AI-transcribed by Claude Opus 5.5 (claude-opus-5-5, adaptive thinking, effort high), Message Batches API, batch ega24b (msgbatch_01RGJJeHdR96U16ZWFuUjtDx), 2026-11-13; not yet checked by eye.
\scanpage{24/29}
\mspage{29}

Let $U_{P}$ be the set of $x \in X_{P}$ such that if $\xi$ denotes the image of $x$
in $P^{\vee}$ the property P should be satisfied by $G_{\xi}$ (resp. $Y_{\xi}$) at the
point $x$ and such that $\o_{\xi}$ is $F_{K(\xi)}$ regular at $x$.  Then $U_{P}$ is open and
$G/U_{P}$ is flat relative to S.

Indeed by the very definition we have $U_{P}$ \quad $U$ (notation of 5.3
(a)) and we apply Par. 12 to $U \to P^{\vee}$ and $F_{\underline{P}^{\vee}}/U$.

\uline{Corollary 5.7}

Let us suppose that F is of finite presentation flat relative
to S, and supp F proper over S (e.g. X proper over S).  Let $V_{P}$ be
the set of $\xi \in P^{\vee}$ such that $G_{\xi}$ (resp. Y) satisfies the property P
and that is $F_{K(\xi)}$ regular.  Under these conditions $V_{P}$ is open, (and
it is \uline{anyway} constructible in every case i.e. without any
assumption of flatness or of properness).

It seems to me that from the point of view of presentation we
cannot leave 5.6 as is with a simple reference to another chapter
[\uline{fascicule}], but it requires an explicit list (i), (ii),... of
properties which we have in view.  Also remark (in 5.1 perhaps)
that the case P = geometrically normal (with S = Spec(k) for sure)
is due to Seidenberg.

\clearpage
% EGA V, English translation by P. Blass, 1987 (cote 24): archive cote 24, site page 30 (PDF page 31).
% typescript. Draft compared: none.
% AI-transcribed by Claude Opus 5.5 (claude-opus-5-5, adaptive thinking, effort high), Message Batches API, batch ega24b (msgbatch_01RGJJeHdR96U16ZWFuUjtDx), 2026-11-13; not yet checked by eye.
\scanpage{24/30}
\mspage{30}

\uline{Section 6}

\uline{Connectedness of an arbitrary hyperplane section}.

We shall here combine the already known criterion of
geometric connectedness of the generic hyperplane section (3.3)
with the connectedness theorem of Zariski in order to obtain a
connectedness result for an arbitrary hyperplane section:

\uline{Proposition 6.1}

We suppose $S = \mathrm{Spec}(k)$ \ , \ $k$ \ an algebraically closed
field [x $X$ proper over $k$ suppose]* that for every irreducible
component $X_i$ of $X$, $f(X_i)$ should be of dimension $\geq 2$, finally that $X$
cannot be disconnected by a closed subset $Z$ of $X$ such that dim
$f(Z) \leq 0$ under such conditions for every $\xi \in \mathrm{Pv}$ \quad $\underline{Y}_{\xi}$ is
geometrically connected.

\uline{Proof}

Since none of the $f(X_i)$ is finite we see that every
irreducible component \ $Y_i$ of $Y$ dominates Pv; on the other hand,
$Y \longrightarrow \mathrm{Pv}$ is proper (if $Y$ is proper over $k$, being such over $X$ which
is proper over $k$). \ On the other hand by (3.3) the generic fiber $Y_{\eta}$
of $\underline{Y} \longrightarrow \mathrm{Pv}$ is geometrically connected.

Finally Pv is regular and à fortiori geometrically unibranch.
It now suffices to apply 15.6.3 \ (which is variant of the Zariski
connectedness theorem) to conclude that \uline{all} the fibers of $\underline{Y} \longrightarrow \mathrm{Pv}$
are geometrically connected \uline{q. ed.} .

Indeed it is not difficult by a proof of analogous type to
generalize 6.1 in the same sense as in 4.4. \ If you don't want to

\clearpage
% EGA V, English translation by P. Blass, 1987 (cote 24): archive cote 24, site page 31 (PDF page 32).
% typescript. Draft compared: none.
% AI-transcribed by Claude Opus 5.5 (claude-opus-5-5, adaptive thinking, effort high), Message Batches API, batch ega24b (msgbatch_01RGJJeHdR96U16ZWFuUjtDx), 2026-11-13; not yet checked by eye.
\scanpage{24/31}
\mspage{31}

trouble yourself with this exercise at least mention it as a
remark.  To say also that we do not discriminate in 6.1 with
regard to hyperplane sections that have an excessive (extra)
dimension.
(From the planning point of view) it might be clearer to group
together all the connectedness questions (including 3.3 and 4.4)
in the same No. (or section).

\uline{Section 7}

\uline{Application to the construction of hyperplane sections and}
\uline{multisections of specified type.}
\uline{7.1}

Let us notice that if S = Spec(k) where  k  is an infinite
field then every non-empty open subset of $P^{v}$ contains a k-rational
point therefore in the notations of 4.1 if E (defined in terms of
a constructible property P) contains the generic point $\eta$ , it
contains a k-rational point therefore there exists a hyperplane
section of X (defined over k) having the property P. On the other
hand, S being again arbitrary it is immediate that for every $s\in S$
and for every point $\xi$ of the fiber $Pv_{s}$ rational over k(s) there
exists a section $\xi$ of $Pv$ on an open neighborhood U of s which
passes through $\xi_{0}$.  If now E is again defined as in 4.1 in terms of
a constructible property P and if we have (for example due to No.
5) the fact that E is open then if $\xi_{0} \in E$ then the section \S{} is a
section of E over U at least if we sufficiently shrink or diminish
U.  Therefore we may construct a hyperplane section $\underline{Y}_{\xi}$ of X over an
open neighborhood U of s such that for every $t \in U$ its fiber $Y_{\xi(t)}$

\clearpage
% EGA V, English translation by P. Blass, 1987 (cote 24): archive cote 24, site page 32 (PDF page 33).
% typescript. Draft compared: none.
% AI-transcribed by Claude Opus 5.5 (claude-opus-5-5, adaptive thinking, effort high), Message Batches API, batch ega24b (msgbatch_01RGJJeHdR96U16ZWFuUjtDx), 2026-11-13; not yet checked by eye.
\scanpage{24/32}
\mspage{32}

at t satisfies the property P . If we do not have à priori $\xi_0$ but
if k(s) is infinite we may combine the two preceding remarks to
obtain a hyperplane section over an open neighborhod of s having
the preceding property. Finally using No. 5 we have a criterion
allowing us to assert that (X resp. F being assumed flat over S
which allows us to apply loc. cit.) $Y_\xi$ resp. $G_\xi$ is also flat over
S. We may therefore replacing X by $Y_\xi$ iterate the previous
construction which allows for example under certain conditions to
construct closer and closer (by successive approximations ???)* a
``multisection'' of $S^1$ of X over an open neighbourhood U of the given
point s, such that $S^1 \to U$ should be finite, flat and with fibers
satisfying the property $\underline{P}$. If k(s) is finite we may be con
shouned to do a change of an étale and surjective base change $S^1 \leftrightarrow$
$\dashrightarrow$ U (U an open neighbourhood of s) before being able to apply
the preceding constructions; indeed under the conditions from the
start of 6.1, if k is finite there does not necessarily exist a
rational point over k in the open non-empty set U, but there
certainly exists a closed point of U, thus a point with values in
a finite extension $k^1$ (necessarily separable) of k ; when k =
k(s), therefore we may, after making a suitable finite étale
extension $S^1$ over a neighbourhood U of s, corresponding to the
residual extension $k^1$ i.e. such that $S^1_s \xrightarrow{\sim}$ Spec($k^1$) restrict
ourselves to the favorable situation of the unique point $s^1 \in S^1$
over s after a base change $S^1 \to S$. I must however note or point
out [un remords Fr] to 4.2 and 4.3 which should have been
announced in a slightly more general form should have been
enumerated at least as a remark : If we are given an integer m

\margin{* Translator's note.}

\clearpage
% EGA V, English translation by P. Blass, 1987 (cote 24): archive cote 24, site page 33 (PDF page 34).
% typescript. Draft compared: none.
% AI-transcribed by Claude Opus 5.5 (claude-opus-5-5, adaptive thinking, effort high), Message Batches API, batch ega24b (msgbatch_01RGJJeHdR96U16ZWFuUjtDx), 2026-11-13; not yet checked by eye.
\scanpage{24/33}
\mspage{33}

and if we denote by E the set of $\xi \in \overset{\vee}{\underline{P}}$ such that $G_\xi$ resp. $\underline{Y}_\xi$
satisfies $\underline{P}$ except over a closed set of dimension $\leq$ m (i.e. the
set $\underline{P}$-singular Z is of dimension $\leq$ m).  Then a) E is a
constructible subset of $\underline{P}\mathrm{v}$ and b) in the case S = Spec(k) if F
respectively X satisfies $\underline{P}$ except over a set of dimension $\leq$ m + 1
then E contains a non-empty open set.

\uline{Proposition 7.2}

Let us assume that X is proper over S and that F is of finite
presentation finite and flat over S.  Let $\underline{P}$ be one of the
properties (i) to (v) of (4.2) and let m be an integer.  Let s $\in$ S
and let us suppose that the set $Z_s$ of points of $X_s$ where $F_s$ does
not satisfy $\underline{P}$ is of dimension $\leq$ m + 1.  Then if also k(s) is
infinite there exists a neighbourhood U of s in E and a section $\xi$
of $\underline{P}\mathrm{v}$ over U having the following properties :  For every s$\in$ U the
set of points of $Y_{\xi,\ s}$ where $G_{\xi(s)}$ does not satisfy $\underline{P}$ is of
dimension $\leq$ m and $\phi_{\xi(s)}$ is $F_{\xi,s}$ regular.  Under such conditions the
module $G_\xi$ over $Y_\xi$ is flat relative to U.  Finally if k(s) is not
supposed infinite we can again make the previous construction
after an étale extension of the type anticipated in 7.1.

\uline{Proposition 7.3}.  Essentially the same.  There is no longer an F
and assume that X is flat relative to S we refer to properties (i)
to (v) of 4.3 in place of those of 4.2 (but being careful to make
the reservation.")  k(s) of characteristic zero or f: X $\rightarrow$ $\underline{P}$ is an

\clearpage
% EGA V, English translation by P. Blass, 1987 (cote 24): archive cote 24, site page 34 (PDF page 35).
% typescript. Draft compared: none.
% AI-transcribed by Claude Opus 5.5 (claude-opus-5-5, adaptive thinking, effort high), Message Batches API, batch ega24b (msgbatch_01RGJJeHdR96U16ZWFuUjtDx), 2026-11-13; not yet checked by eye.
\scanpage{24/34}
\mspage{34}

immersion and in the case (v) that for every s $\in$ S [illegible]
irreducible component Z of X$_s$ we have dim f(Z) $\geq$ 2.
[Nota Bene:  For (v) [[illegible]] 12.2.1   (x) and (xi) (we can
then \uline{illegible} in the other case 4.3 or 12.2.1 (x)] (marginal
remark largely illegible in preceding square brackets).

Text crossed out

\uline{Proposition 7.4}

Let g:   X $\to$ S be a flat proper morphism, let s $\in$ S let us
put n = dim X$_s$ and let us suppose that the dimension of the set of
points of X$_s$ where X$_s$ is not separable over k(s) is  $<$ n.  (\uline{for
example}  X$_s$   separable).  Then there exists an open neighborhood
U of s and an \'etale finite, surjective morphism S1---> U such that
X \quad x$_S$ \quad S1 \quad admits a section over S1 .  If k(s) is infinite we
may take for S$^1$ a closed subscheme of \uline{X\ \ U}.*  Let us assume to
start with that k(s) is infinite.

We proceed by induction on n, the case n = 0 being trivial.
Indeed in that case there exists an open neighbourhood U of s such
that X|U itself is \'etale, finite and surjective above U as we see
by immediate cross references.  If n $>$ 0, we apply 7.3 for the
``separable'' property which allows us to replace X by a ``hyperplane
section'' $\underline{Y}$ having the same properties up to this that n is
replaced by n-1.  If k(s) is not assumed infinite we begin by
making an \'etale base change, it works.

\margin{* unclear, ask a.g.}

\clearpage
% EGA V, English translation by P. Blass, 1987 (cote 24): archive cote 24, site page 35 (PDF page 36).
% typescript. Draft compared: none.
% AI-transcribed by Claude Opus 5.5 (claude-opus-5-5, adaptive thinking, effort high), Message Batches API, batch ega24b (msgbatch_01RGJJeHdR96U16ZWFuUjtDx), 2026-11-13; not yet checked by eye.
\scanpage{24/35}
\mspage{35}

\uline{Remark 7.5}

In particular if $X$ is projective and separable over $S$ it
admits locally over $S$ étale multisections.  But we note that we
can give examples with $X$ proper and smooth (but non-projective) $S$,
where the same conclusion fails.  Of course, the projective
assumption cannot be weakened in general to an assumption of
quasi-projectiveness as we see for example by taking $X$ étale non-%
finite over $S$ ......

ILLEGIBLE *

\clearpage
% EGA V, English translation by P. Blass, 1987 (cote 24): archive cote 24, site page 36 (PDF page 37).
% typescript. Draft compared: none.
% AI-transcribed by Claude Opus 5.5 (claude-opus-5-5, adaptive thinking, effort high), Message Batches API, batch ega24b (msgbatch_01RGJJeHdR96U16ZWFuUjtDx), 2026-11-13; not yet checked by eye.
\scanpage{24/36}
\mspage{36}

\uline{Section 8}

\uline{Dimension of the set of exceptional hyperplanes}

\uline{8.1} In the previous sections and notably sections 2 and 3 we have
given statements asserting that the set of $\xi \in \underline{P}v$ such that the set
of $\xi \in \underline{P}v$ such that $\underline{Y}$ has a certain property $\underline{P}$ is constructible and
that it contains the generic point ${}_{\eta}$ or else that the set $Z_{\underline{P}}$ of $\xi \in$
$\underline{P}v$ ``exceptional for $\underline{P}$'' is constructible and is rare i.e. that its
closure is of codimension $\geq 1$.  (Nota Bene:  we suppose that $S =$
$\mathrm{Spec}(k)$).

In certain cases we can make this statement more precise by
giving a better upper bound for this codimension, which is
important for certain questions.  For example, if we see that this
codimension is greater than or equal to two it follows that a
``sufficiently general'' straight line $D$ of $\underline{P}v$ does not intersect $Z_{\underline{P}}$,
whence the existence (if $k$ is infinite) of ``linear pencils'' of
hyperplane sections $Y_{\xi}$ ($\xi$ - a geometric point of $D$) all of which
have the property $\underline{P}$ (see section number [omitted ask
Grothendieck]* for examples).

From the \uline{writing up} point of view, since the results of the
present N0 make more precise some results of the previous sections
the question arises if it is necessary to do this catching up in a
separate section (or number) or to give a more precise version
gradually as we move along.  Redactor decidetur (Latin).

\uline{8.2} Let $Z$ be the set of $\xi \in \underline{P}v$ such that $\dim Y_{\xi} > \dim X - 1$ and
let us suppose that for every irreducible component irr $X_i$ of $X$ we

\clearpage
% EGA V, English translation by P. Blass, 1987 (cote 24): archive cote 24, site page 37 (PDF page 38).
% typescript. Draft compared: none.
% AI-transcribed by Claude Opus 5.5 (claude-opus-5-5, adaptive thinking, effort high), Message Batches API, batch pilot (msgbatch_0189Nhf1xB615YWnztubtLfc), 2026-11-12; not yet checked by eye.
\scanpage{24/37}
\mspage{37}

have dim f($X_{i}$) > [illegible] [is it two ask Grothendieck]* then Z
is of codimension two in $\underline{P}$v.  This follows from 2.1 and 2.2.
(which implies that every irreducible component of $\underline{Y}$ dominates $\underline{P}$v)
and from the dimension theory for the morphism $\underline{Y} \rightarrow \underline{P}$v.  Starting
from this result we may give as a corollary the case where we
start from a closed subset Z of X and where we consider the
dimension of the inverse images $Z_{\xi}$ in the $Y_{\xi}$ ($\xi \in \underline{P}$v) and we may
even take for Z the set of $\xi \in \underline{P}$v such that there exists an
irreducible component of $T_{k(\xi)}$ whose trace on $Y_{\xi}$ has the greatest
dimension (NB we always assume that for every irreducible $Z_{i}$ of Z
we have dim f($Z_{i}$) > 0).

Finally the most precise statement in this direction and one
that results easily from the first announcement (for X irreducible) and
from 2.7 is the following modified statement:  F being coherent
over X, suppose that for every associated prime cycle T for F we
have dim f(T) > 0 then the set of $\xi \in \underline{P}$v such that $\phi_{\xi}$ is not $F_{k(\xi)}$ -
regular is (constructible and) of codimension $\geq$ 2.  (The notation
for $\phi_{\xi}$ is that from N05).  We can give this as the principal
assertion, and announce the previous assertions as corollaries,
the proof being or proceeding via one of the corollaries.

Please note that with the preceding notations if $\xi \in \overset{\mathrm{v}}{\underline{P}}$ - Z,
then for every y $\in \underline{Y}_{\xi}$ we have coprof$_{y}$ $G_{k(\xi)}$ = coprof$_{y}$ $G_{\xi}$ and
consequently if coprof F $\leq$ n in particular if F is Cohen-Macauley
then for $\xi \in \underline{P}$v - Z, $G_{\xi}$ is Cohen-Macauley.  Finally if F is ($S_{k}$)
[illegible ask Grothendieck]*

$S_{k-1}$ for $\xi \in \underline{P}$v - Z (reference $O_{IV}$).

\margin{* Ask AG.}

\clearpage
% EGA V, English translation by P. Blass, 1987 (cote 24): archive cote 24, site page 38 (PDF page 39).
% typescript. Draft compared: none.
% AI-transcribed by Claude Opus 5.5 (claude-opus-5-5, adaptive thinking, effort high), Message Batches API, batch ega24b (msgbatch_01RGJJeHdR96U16ZWFuUjtDx), 2026-11-13; not yet checked by eye.
\scanpage{24/38}
\mspage{38}

\uline{8.3}

We notice that for [illegible ask Grothendieck]* even if
$k = 1$, $X$ being geometrically integral of dimension two where ($k =$
2 ???) $X$ being geometrically integral and geometrically normal of
dim 3).  It is enough to start from a projective integral surface
\[ X \ C \ \underline{P}v \]
over $k$ algebraically closed having a point $x$ where $X$ is not Cohen-%
Macauley, then for every hyperplane pressing through $x$ the
corresponding hyperplane section $Y_{\xi}$ admits $x$ as an associated
embedded cycle (respectively we start from a normal (thus $S_2$)
integral variety $X \ C\underline{P}v$ of dimension three having a point $x \in X$
where $X$ is not Cohen-Macauley, then the $Y$'s passing through $x$ are
not CM i.e. they are met ($S_2$) at $x$.)

In these examples the set of ``exceptional'' $\xi$ for the property
($S_k$) contains the hyperplane of $\underline{P}v$ defined by $x \in \underline{P}$ and it is of
codimension one (and not of codimension $\geq$ two) compare 8.5 below
for a general precise result in this direction along these lines?

Footnote(?)  Fairly illegible ask Grothendieck*

If $F$ is ($S_k$) [illegible]
Illegible Gap  5 lines.

We point out the following result:

\margin{* ask AG}

\clearpage
% EGA V, English translation by P. Blass, 1987 (cote 24): archive cote 24, site page 39 (PDF page 40).
% typescript. Draft compared: none.
% AI-transcribed by Claude Opus 5.5 (claude-opus-5-5, adaptive thinking, effort high), Message Batches API, batch ega24b (msgbatch_01RGJJeHdR96U16ZWFuUjtDx), 2026-11-13; not yet checked by eye.
\scanpage{24/39}
\mspage{39}

\uline{Proposition 8.4}

Let T be a closed subset of X and suppose that codim (T, X) $\geq$
k. Then for every $\xi \in \underline{P}$v we have codim (T$_{\xi}$, $\underline{Y}_{\xi}$) $\geq$ k - 1. Let Z be
the set of $\xi \in \underline{P}$v such that codim (T$_{\xi}$, $\underline{Y}_{\xi}$) = k - 1 (i.e. codim (T$_{\xi}$,
$\underline{Y}_{\xi}$) < k) then Z is a constructible, nowhere dense [rare Fr]
subset of $\overset{\text{v}}{\underline{P}}$, i.e. $\overline{\text{Z}}$ is of codimension $\geq$ \unsure{$|$|1} in $\underline{P}$v.

In order for it to be of codimension $\geq$ 2 it is necessary and
sufficient that for every irreducible component T$_{\text{i}}$ of X of
codimension equal to k and such that dim f(T$_{\text{i}}$) = 0 there should
exist one component \unsure{$|$|1} [illegible is there an insert a footnote ask
Grothendieck].*

The first assertion follows immediately from the following lemma
8.4.1 (a) which is a remorseful afterthought to paragraph 5.

\uline{Lemma 8.4.1}

Let X be a locally noetherian prescheme, let L be an
invertible module over X, $\phi$ a section of L, $\underline{Y}$ = V($\phi$), T a closed
subset of X. Let us assume that codim ($\underline{Y}$, X) $\geq$ K.

Then

a) Codim (T$\cap$ Y, $\underline{Y}$) $\geq$ k - 1

b) In order to have

codim (T$\cap$ Y, Y) = k - 1

i.e. codim (T$\cap$ Y, Y) < k

it is necessary and sufficient that there should exist an
irreducible component T$_{\text{i}}$ of T contained in Y, and such that codim

\margin{$*$ Ask Grothendieck}

\clearpage
% EGA V, English translation by P. Blass, 1987 (cote 24): archive cote 24, site page 40 (PDF page 41).
% typescript. Draft compared: none.
% AI-transcribed by Claude Opus 5.5 (claude-opus-5-5, adaptive thinking, effort high), Message Batches API, batch ega24b (msgbatch_01RGJJeHdR96U16ZWFuUjtDx), 2026-11-13; not yet checked by eye.
\scanpage{24/40}
\mspage{40}

$(T_i, X) = k$ and such that for every irreducible component $X_j$ of $X$
containing $T_i$ and such that
\[ \dim O_{X_j, T_i} = \dim O_{X, T_i} \]
i ( = k) we have
\[ X_j \over{C}{\not\subset} Y. \]

The verification of this lemma is immediate due to the general
facts in $O_{IV}$ chapter IV about dimension.  With the assumptions of
8.4, by 8.4.1(b) we see which ones are the exceptional hyperplanes
$H_\xi$.  If we exclude the set $Z_0$ of $\xi \in \underline{P}\mathrm{v}$ such that there is an
irreducible component R of T or of X such that dim f(R) > 0 and
such that $R_\xi$ is of ``dimension too large'' (a set which is of
codimension two and if what follows it does not count the
exceptional $H_\xi$ are there for which there exists a $T_i$ with codim
$(T_i, X) = k$ and dim $f(T) = 0$, $f(T)$ CH and such that for \uline{every}
irreducible component $X_j \supset T_i$ of X with codim $(T_i, X_j) = k$ we have
$f(X_j) \not\subset H_\xi$.  For a given $T_i$ with codim $(T_i, X) = k$ if there exists
an $X_j$ with codim $(T_i, X_j)$ = ? [illegible ask Grothendieck]
and such that dim $f(X_j) = 0$ then we will have
\[ f(X_j) = f(T_i)\ \mathrm{C}\ H_\xi \text{ and consequently } \xi \text{ would not be exceptional} \]
relative to the $T_i$ if and only if $f(T_i)$ C $H_\xi$ ; the set of such $\xi$ is
(the trace over) of a hyperplane of $\underline{P}\mathrm{v}$.  This proves
$\qquad\qquad P - Z_0$

8.4, and also proves the more precise result that the exceptional
set is the union of a set of codimension $\geq 2$ and of a \uline{union of
hyperplanes} determined in an evident way by the above proof.

\clearpage
% EGA V, English translation by P. Blass, 1987 (cote 24): archive cote 24, site page 41 (PDF page 42).
% typescript. Draft compared: none.
% AI-transcribed by Claude Opus 5.5 (claude-opus-5-5, adaptive thinking, effort high), Message Batches API, batch ega24b (msgbatch_01RGJJeHdR96U16ZWFuUjtDx), 2026-11-13; not yet checked by eye.
\scanpage{24/41}
\mspage{41}

(I am afraid that the write up is quite floppy or perhaps
sloppy since I have reasoned geometrically all the time without
saying so, by taking points over an algebraically closed field.
Of course, the condition announced in 8.4 is indeed geometric so
that we may suppose k alg-closed and argue for K-rational points)
using 8.4; 5.7.4 and the end of 8.2 we obtain:

\uline{Corollary 8.5}

Suppose that for all associated prime cycles we have at most
simply [illegible]* and suppose that F satisfies $(S_k)$.

The (constructible) set of points of $\underline{P}^{\vee}$ such that $\phi_\xi$ is $F_{k(\xi)}$
regular and $G_\xi$ is $(S_k)$ has a complement of codimension at least two
if and only if ($\Leftrightarrow$) for every integer $n \geq 0$ we denote by $Z_n$ the set
of $x \in T = \operatorname{supp} F$ such that the $\operatorname{coprof}_x$ [illegible]*

% NOTE: blank space of several lines left on the page here

we see that for every irreducible component $Z_{ni}$ of $Z_n$ with codim
$(Z_{ni},\ T) = n + k + 1$ and $\dim f(Z_{ni}) = 0$ there exists an irreducible
component $T_j$ of $T$ containing $Z_{ni}$ and such that codim $(Z_{ni},\ T_j) = n +$
$k + 1$ and $\dim f(T_j) = 0$.

When f is given finite then for every closed subset R of
[illegible ask A.G] we have $\dim f(R) = \dim R$ so that the criterion
takes the following form: there does not exist an isolated point
z in any one of the $Z_n$ such that $\dim_z T( = \dim F_z)$ is equal to $n +$
$k + 1$.

\margin{* Ask A.G}

\clearpage
% EGA V, English translation by P. Blass, 1987 (cote 24): archive cote 24, site page 42 (PDF page 43).
% typescript. Draft compared: none.
% AI-transcribed by Claude Opus 5.5 (claude-opus-5-5, adaptive thinking, effort high), Message Batches API, batch ega24b (msgbatch_01RGJJeHdR96U16ZWFuUjtDx), 2026-11-13; not yet checked by eye.
\scanpage{24/42}
\mspage{42}

When F is equidimensional of dimension d this condition is vacuous
if $d \leq k$ (and indeed we knew it because in this case the
[hypothesis] ($S_k$) on F is nothing else but the hypothesis CM), and
if $d \geq k + 1$ it means that the set $Z_{d-(K + 1)}$ of points of $T$ where
the co-depth of F is $> d - (k + 1)$ i.e. true depth of $F \geq k + 1$
(even though, by preference we only have true depth of $F \geq k$ as a
consequence of the property ($S_k$) and $k \leq d$). If we no longer
assume that F is equidimensional there remains that we may express
the desired condition in the following simple way:

\uline{8.6} For every closed point $x \in \operatorname{supp} F$ such that $\dim F_x \geq k + 1$ we
have $\operatorname{prof} F_x \geq k + 1$

The sufficiency is seen immediately by putting $Z = X$. The
necessity is seen by noticing that for every $\xi$ such that $\phi_\xi$ is $X_{k(\xi)}$
- regular and $x \in Y_\xi$ we have $\dim G_\xi$, $x = \dim F_x - 1$, $\operatorname{prof} G_{\xi\; x} =$
$\operatorname{prof} F_x - 1$ so that x put by default the above condition we have
$\operatorname{prof} G_\xi\ x \geq k$ but $\dim G_\xi\ x \geq k$ which shows that $G_\xi$ does not satisfy
condition ($S_k$) at x; but the set of $\xi$ such that $x \in Y_\xi$ is of
codimension 1 (NB: I implicitly assumed that k is algebraically
closed the case to which we reduce immediately). The preceding
general criterion should be evident in the case 8.6. We now study
the points y of $\underline{Y}$ that are not smooth for $\underline{Y}_\xi$ relative to $k(\xi)$. We
restrict ourselves to the case where $f\colon X \to f$ is \uline{unramified}
(practically, it will be an immersion). and where $X \to S$ is
smooth. We do not necessarily assume that S is the spectrum of a
field. Since f is unramified the canonical homomorphism $f^*(\Omega^1{}_{\underline{P}/S})$
$\to \Omega^1{}_{X/S}$ is surjective and its kernel is a locally free module

\margin{* unclear, ask AG.}

\clearpage
% EGA V, English translation by P. Blass, 1987 (cote 24): archive cote 24, site page 43 (PDF page 44).
% typescript. Draft compared: none.
% AI-transcribed by Claude Opus 5.5 (claude-opus-5-5, adaptive thinking, effort high), Message Batches API, batch ega24b (msgbatch_01RGJJeHdR96U16ZWFuUjtDx), 2026-11-13; not yet checked by eye.
\scanpage{24/43}
\mspage{43}

over X which we denote $\hand{\unsure{\mathcal{N}}}^{\mathrm{v}}{}_{X/\underline{P}}$; when if f is an immersion this is
nothing else but the conormal module $J/J^2$ defined by the ideal J
of X in $\underline{P}$ and we call it in every case the conormal module.
\[ \text{(a)}\quad 0 \rightarrow \hand{\unsure{\mathcal{N}}}^{\mathrm{v}}{}_{X/\underline{P}} \rightarrow f^*(\Omega^1{}_{\underline{P}/S}) \rightarrow \Omega^1{}_{X/S} \rightarrow 0 \]
Let us observe that we have also over $\underline{P}$ an exact canonical
sequence (which should appear as an example in paragraph 16 for
example)
\[ \text{(b)}\quad 0 \rightarrow \Omega^1{}_{\underline{P}/S}(1) \rightarrow E_{\underline{P}} \rightarrow O_{\underline{P}}(1) \rightarrow 0 \]
(i.e. $\Omega^1{}_{\underline{P}/S}$ is canonically isomorphic to the kernel of the
canonical homomorphism) $E_{\underline{P}}(-1) \rightarrow O_{\underline{P}}$ deduced from
$E_{\underline{P}} \rightarrow O_{\underline{P}}(1)$, to it we apply $f^*$:
\[ \text{(b}^1\text{)}\quad 0 \rightarrow f^*(\Omega^1{}_{\underline{P}/\underline{S}}\hand{\unsure{)}}(1) \rightarrow E_X \rightarrow O_X(1) \rightarrow 0 \]
which gives an explicit description of $f^*(\Omega^1{}_{\underline{P}/\underline{S}}\hand{\unsure{)}}(1)$ over X and
allows therefore to identify $\hand{\unsure{\mathcal{N}}}^{\mathrm{v}}{}_{X/\underline{P}}(1)$ with a submodule locally a
direct factor of $E_X$ or again the dual $\hand{\unsure{\mathcal{N}}}{}_{X/\underline{P}}(-1)$ is canonically
isomorphic to a quotient module of $E^{\mathrm{V}}{}_X$. Consequently $\underline{P}(\hand{\unsure{\mathcal{N}}}{}_{X/\underline{P}}(-1)$
$= \underline{P}(\hand{\unsure{\mathcal{N}}}{}_{X/\underline{P}})$ can be canonically embedded into $\underline{P}(E^{\mathrm{V}}{}_X) = X\mathrm{x}_S\underline{P}^{\mathrm{v}} = X_{\underline{P}}{}^{\mathrm{V}}$
as a projective sub-fibration over X therefore as a closed
subscheme. The latter is necessarily contained in $\underline{Y}$ (from the
fact that $\Omega_{X/\underline{P}}(1)$ is contained in the kernel of $E_x \rightarrow O_X(1)$ [last
two symbols illegible ask A G]

The underlying set of this prescheme is nothing else but the
set of points of $\underline{Y} = V(\emptyset)$ which are \uline{singular zeros} (par. 16) of
the section $\emptyset$ of $O_X\mathrm{x}_{\underline{P}}\mathrm{v}(1, 1)$ relative to the base $\underline{P}^{\mathrm{v}}$ i.e. its
points with values in the field k over $\underline{P}^{\mathrm{v}}$ are the points x of
$Y_k \hand{\subset} X_k$ such that $\emptyset_k$ vanishes to order at least two at x i.e.
such that $Y_k$ is not smooth of relative dimension r - 1 over k at

\clearpage
% EGA V, English translation by P. Blass, 1987 (cote 24): archive cote 24, site page 44 (PDF page 45).
% typescript. Draft compared: none.
% AI-transcribed by Claude Opus 5.5 (claude-opus-5-5, adaptive thinking, effort high), Message Batches API, batch ega24b (msgbatch_01RGJJeHdR96U16ZWFuUjtDx), 2026-11-13; not yet checked by eye.
\scanpage{24/44}
\mspage{44}

x.  The announced characterization of singular zeros [illegible
ask AG] the elements of a smooth subscheme $\underline{P}(\hand{\unsure{\nu}}\ X/\underline{P})$ of $X_{\underline{P}}{}^{v}$ gives
immediately the following statement which deserves to appear as a
preliminary proposition if S = Spec k and if H is a hyperplane of
$\underline{P}$ then $Y = X \ \mathrm{x}_{\underline{P}}\ H$ is smooth over k of relative dimension (d - 1)
at the point $x \in \underline{Y}(k)$ (i.e. x is a non-singular zero i.e.
geometrically non-singular of the section {\o} of $O_X(1)$ defined by H)
if and only if H does not contain the image by {\o} of the tangent
space to X at x (relative to k) or as we say once more (if
$f{:}X \to \underline{P}$ is an immersion which allows us to identify X to a
subscheme of $\underline{P}$) if and only if H is not tangent to X at x.  This
follows trivially from the Jacobian criterion of smoothness or
from the definition of a singular zero, once we make precise the
sense of the statement, that is to say that we make precise how a
vector subspace of the tangent space to $\underline{P}$ at a point a(= f(x))
defines a linear subspace of $\underline{P}$ (in such a way that it makes sense
to say that H does not contain the soul vector** subspace): of
course this comes from the exact sequence (b) above which allows
to define a one-to-one correspondence between the set of vector
subspaces of the tangent space at a and the set of linear
subspaces to $\underline{P}$ (?)* [illegible ask AG] containing a.  This
correspondence anyway reduces to associating to a linear
subvariety passing through a its tangent space at a considered as
a subspace of the tangent space to $\underline{P}$ at a.

Such ``sorites'' grouped together with various ``sorites'' about
linear subvarieties and about grossmanians ought to be given in
one or to preliminary numbers or paragraphs of course announcing

\margin{** Check this Blass}

\clearpage
% EGA V, English translation by P. Blass, 1987 (cote 24): archive cote 24, site page 45 (PDF page 46).
% typescript. Draft compared: none.
% AI-transcribed by Claude Opus 5.5 (claude-opus-5-5, adaptive thinking, effort high), Message Batches API, batch ega24b (msgbatch_01RGJJeHdR96U16ZWFuUjtDx), 2026-11-13; not yet checked by eye.
\scanpage{24/45}
\mspage{45}

them over any base.  In fact we can do better knowing that the
prescheme $\underline{Y}^{sing}$ of singular zeros of \o{} relative to $\underline{P}^{v}$ defined in
par. 16 is nothing else but $\underline{P}(\hand{\unsure{\mathcal{J}|\mathcal{N}}}{}^{v}{}_{x/P})$ and (since the latter is
smooth over S of relative dimension dt(r - d - 1) = r - 1 (r being
the relative dimension of $\underline{P}^{v}$ over S) we are under the favorable
conditions studied in N$^\circ$ 16 or paragraph 16.*

In order to verify them, let us notice that by definition
$\underline{Y}^{sing}$ is nothing else but the sub-prescheme of $\underline{Y}$ of zeros of the
section $\psi$ = d \o/$\underline{Y}$ of $\Omega^{1}{}_{X_{\underline{P}}{}^{v}/\underline{P}^{v}}\ (1,\ 1) \otimes O_{y} = \Omega^{1}{}_{X/S} \otimes O_{\underline{Y}}(1,\ 1)$
[illegible]**

We shall give another interpretation of this section from
which the conclusion follows immediately.  In order to do this let
us consider the following diagram of exact sequences over $X_{\underline{P}}{}^{v}$ or
more generally over any prescheme Z over $X_{\underline{P}}{}^{v}$.

% NOTE: the rest of the page is blank (no diagram typed)

\margin{* Ask AG about this reference}
\margin{** Ask AG}

\clearpage
% EGA V, English translation by P. Blass, 1987 (cote 24): archive cote 24, site page 46 (PDF page 47).
% typescript. Draft compared: none.
% AI-transcribed by Claude Opus 5.5 (claude-opus-5-5, adaptive thinking, effort high), Message Batches API, batch ega24b (msgbatch_01RGJJeHdR96U16ZWFuUjtDx), 2026-11-13; not yet checked by eye.
\scanpage{24/46}
\mspage{46}

Diagram:

\begin{diagramblock}
\begin{tikzcd}
 & & 0 & & \\
 & 0 & G_{\underline{P}}{}^{v}{}_{/S}\otimes O_Z(0,-1) \arrow[u] & & \\
 & \Omega^1{}_{X/Y}\otimes O_Z(1,0) \arrow[u] & & & \\
0 \arrow[r] & \Omega^1{}_{\underline{P}/S}\otimes O_Z(1,0) \arrow[u, dash] \arrow[r] & E\otimes O_Z \arrow[uu] \arrow[r] & O_Z(1,0) \arrow[r] & 0 \\
 & {}^{v}{}_{x/\underline{P}}\otimes O_Z(1,0) \arrow[u] & O_Z(0,-1) \arrow[u] \arrow[ul, dashed, "B"'] \arrow[ur, dashed, "\partial"] & & \\
 & 0 \arrow[u] & 0 \arrow[u] & &
\end{tikzcd}
\end{diagramblock}

[Note to AG the upper G is really an illegible letter \hand{\unsure{$\mathcal{Y}$|$\mathcal{G}$}} $\underline{P}^{v}/S$
what is this?]*
where the first column is deduced from (a) by tensoring with
$O_Z(1,0)$ the row is deduced from (b) by tensoring with $O_Z$ and the
column two is deduced from its transpose from the analogous
sequence ($b^{v}$) relative to $\underline{P}^{v}$ (obtained by replacing E by $E^{v}$) and
tensoring with $O_Z$.  From the very definition of $\underline{Y}$, Z is over $\underline{Y}$ if
and only if the composed morphism $\partial$ from the diagram is zero i.e.
if we can find a factorization B: $O_Z(0,-1) \rightarrow \Omega^1{}_{\underline{P}/S} \otimes O_Z(1,0)$.  If
this is the case we can consider it composition with $\Omega^1{}_{\underline{P}/S} \otimes$
$O_Z(1,0) \rightarrow \Omega^1{}_{X/Y} \otimes O_Z(1,1)$.  I say that this is precisely the
section $\psi$ [Blass: check if this letter is OK]** which we have

\margin{* Ask AG}
\margin{** Blass check this}

\clearpage
% EGA V, English translation by P. Blass, 1987 (cote 24): archive cote 24, site page 47 (PDF page 48).
% typescript. Draft compared: none.
% AI-transcribed by Claude Opus 5.5 (claude-opus-5-5, adaptive thinking, effort high), Message Batches API, batch ega24b (msgbatch_01RGJJeHdR96U16ZWFuUjtDx), 2026-11-13; not yet checked by eye.
\scanpage{24/47}
\mspage{47}

introduced above (the verification ought to be essentially
mechanical). It is zero if and only if Z is above lies over $V(\psi)$
(by the very definition of $V(\psi)$!) but this means also that B can
be factored by $\hand{\unsure{\mathcal{N}}}^{v}{}_{x/\underline{P}} \otimes O_{Z}(1,0)$ i.e. that the sub-module $O_{Z}(0,-1)$
of $E \otimes O_{Z}$ is contained in the sub-module $\hand{\unsure{\mathcal{N}}}\;{}_{x/\underline{P}} \otimes O_{Z}(1,0)$ which
evidently signifies also that Z is over the sub-prescheme $\underline{P}(\hand{\unsure{\mathcal{N}}}$
${}_{x/\underline{P}}(1))$ of $\underline{P}(E^{v}{}_{X})$, achieving the proof that we have announced just
before this erudite exercise in syntax for which I have already
had to sweat quite a bit we could remark that from every set
theoretic point of view $Y^{sing}$ is of dimension $r - 1$ if $S = \mathrm{Spec}\ k$,
whereas $\underline{P}^{v}$ is of dimension r so that the image of $\underline{Y}^{sing}$ in $\underline{P}^{v}$ is
of codimension $\geq 1$ which gives again 2.12 (it is well to note that
the argument is not essentially distinct from the one used in
2.12). We note that most often this set is effectively of
codimension one (compare below).

Consequently we cannot in general find the ``linear pencils''
of hyperplane sections all of which are smooth. However we shall
see that we can often manage to find the pencils formed by
hyperplane sections not having any supersingular point due to the
fact that in the most common cases the image of $Y^{sup\ sing}$ in $\underline{P}^{v}$ is
of codimension two.

We shall first of all recall the essential points
differential in nature of the situation studied here:

\clearpage
% EGA V, English translation by P. Blass, 1987 (cote 24): archive cote 24, site page 48 (PDF page 49).
% typescript. Draft compared: none.
% AI-transcribed by Claude Opus 5.5 (claude-opus-5-5, adaptive thinking, effort high), Message Batches API, batch ega24b (msgbatch_01RGJJeHdR96U16ZWFuUjtDx), 2026-11-13; not yet checked by eye.
\scanpage{24/48}
\mspage{48}

\uline{Theorem 8.7}

(a) The sub-prescheme $\underline{Y}^{\mathrm{sing}}$ (defined in N$^\circ$ or par. 16) in the
present situation is nothing else but $\underline{P}(\hand{\unsure{\mathcal{N}}}_{\ x/\underline{P}})$ considered as a
sub-scheme of $\underline{Y}$ as explained above.

(b) The underlying set of the prescheme $\underline{Y}^{\mathrm{sup\ sing}}$ (cf N$^\circ$ or
par. 16) is nothing else but the set of ramification points of
morphism of smooth preschemes over S of relative dimension r - 1
and r (namely) $\underline{Y}^{\mathrm{sing}} = \underline{P}(\hand{\unsure{\mathcal{N}}}_{\ x/\underline{P}}) \rightarrow \underline{P}^{\mathrm{v}}$ i.e. in order for the latter
morphism to be unramified at the point y (ref to the definition)
it is necessary and sufficient that y should be geometrically an
ordinary singular point for $\text{ø}_{\hand{\xi}}$ ($\hand{\xi}$ being the point of $\underline{P}^{\mathrm{v}}$ that is the
image of y).

(c) Let us assume S = Spec(k) and that y $\in$ $Y^{\mathrm{sing}} = \underline{P}(\hand{\unsure{\mathcal{N}}})$ is
a k - rational point, let [illegible]* and $\hand{\xi}$ be its projections in
X(k) resp $\underline{P}^{\mathrm{v}}$(k) and let us consider the linear subvariety H$^1$ of $\underline{P}^{\mathrm{v}}$
``image'' of the tangent map of the closure of its [Fr. son image]
image in $\underline{P}^{\mathrm{v}}$, given the induced reduced structure and let us
consider the induced morphism g:$\underline{Y}^{\mathrm{sing}} \rightarrow$ T (a dominant morphism of
integral preschemes). The conditions (i) and (ii) (bis) are
equivalent:

(i) The morphism g is generically étale (i.e. étale at least
one point or what is the same is étale = unramified at the generic
point of $\underline{Y}^{\mathrm{sing}}$)

(i bis) The field extension L/K defined by g is finite and
separable

(i ter) The morphism g is birational i.e. the extension L/K
is the trivial extension.

\margin{* Ask AG}

\clearpage
% EGA V, English translation by P. Blass, 1987 (cote 24): archive cote 24, site page 49 (PDF page 50).
% typescript. Draft compared: none.
% AI-transcribed by Claude Opus 5.5 (claude-opus-5-5, adaptive thinking, effort high), Message Batches API, batch ega24b (msgbatch_01RGJJeHdR96U16ZWFuUjtDx), 2026-11-13; not yet checked by eye.
\scanpage{24/49}
\mspage{49}

[Marginal note illegible ask A.G]\hand{*}

(ii) $Y$sing $\neq$ $\underline{Y}$sup sing (set theoretically speaking let us say
(ii ter) There exists an $x \in X(K)$ and a tangent hyperplane H to $X_{K}$
at x which is not osculating at x by which we understand precisely
that x is not supersingular for the section of $O_{X_{K}}$ (1) that defines
H...)

These conditions imply that $\underline{Y}$sup sing $\neq \phi$ [Fr illegible ask
A.G.] dim $\underline{Y}$sup sing $\leq r - 2$ so that the image of Tsup sing in $\underline{P}$v has a
codimension $\geq 2$, and they imply also (iii) dim $T = r - 1$, i.e. T
is of codimension one in $\underline{P}$v.

\uline{Proof}

The equivalence of (i) and (i leir) is trivial its
equivalence with (ii) is a trivial consequence of 8.7 b), finally
the equivalence of (ii) and of (ii ter) is practically the
definition of $\underline{Y}$sup sing. Evidently (i ter) $\Rightarrow$ (i) it remains to
prove that (i) $\Rightarrow$ (i ter). We may evidently suppose that K is
algebraically closed and we are reduced to prove (taking into
account the hypothesis (i)) that there exists an open set $U \neq \phi$
such that $\xi \in U$ (K) implies that there exists exactly one point of
$\underline{Y}$sing (K) over $\xi$. This will follow from 8.7 c) which implies more
precisely.

\uline{Corollary 8.9}

Suppose that condition (i) of 8.8 is satisfied and let U be
the open subset of T of the points where T is smooth over K. Then

\margin{* Ask AG}

\clearpage
% EGA V, English translation by P. Blass, 1987 (cote 24): archive cote 24, site page 50 (PDF page 51).
% typescript. Draft compared: none.
% AI-transcribed by Claude Opus 5.5 (claude-opus-5-5, adaptive thinking, effort high), Message Batches API, batch ega24b (msgbatch_01RGJJeHdR96U16ZWFuUjtDx), 2026-11-13; not yet checked by eye.
\scanpage{24/50}
\mspage{50}

$U \neq 0$ Ysing / U $\Rightarrow$ U is an open immersion a fortiori $\underline{Y}$sing / U does \margin{?}
not contain the points of Ysupsing.  If X is proper over K, then
g : $\underline{Y}$sing $\to$ T is surjective thus $\underline{Y}$sing / U $\Rightarrow$ U is proper over K so
that g : $\underline{Y}$sing $\to$ T is surjective therefore $\underline{Y}$sing / U $\Rightarrow$ U is an
isomorphism and U is the biggest open set of T having the latter \margin{?}
property.

First of all since g is dominating and generically étale it
is generically étale at least one non empty open subset V of T
such that $\underline{Y}$sing $|$ V $\Rightarrow$ V is étale and surjective which implies that
V is smooth over K.  If then $\xi \in$ V(K) and if y is a point of $\underline{Y}$sing
(K) over $\xi$ then with the notations of 8.7 c) the space H$^1$ is
nothing else but the tangent space to T and $\xi$, and as we observed
here this implies that the point x of X(K), the projection of y is
determined as the orthogonal point to H$^1$ thus it is uniquely
determined thus since Ysing $\langle$ X x $\underline{P}$v, y is uniquely determined.

This proves already that g is birational (being generically
étale and generically radicial).  On the other hand the morphism $\psi$
(whose definition in its form is evident) which associates to
every $\xi \in$ U(K) the unique point x  $= \psi(\xi) \in \underline{P}$ orthogonal to the
tangent space to U at $\xi$, coinciding over V with the composition V
$\to \underline{Y}$sing $|$ V $\to$ X, where the second arrow is the projection;
therefore setting h = $(\psi$, id):U $\to \underline{P}$ x $\underline{T}$ (illegible)*
g$_1$ = g $|$ g$^{-1}$(U) : g$^{-1}$(U) $\to$ U the composition hg$_1$ : g$^{-1}$(U) $\to \underline{P}$ x T
is nothing else but the canonical inclusion, this being so for its
restriction to g$^{-1}$(V) $\xrightarrow{\sim}$ V.  It follows that h
factors through the scheme theoretic closure Y$^1$ of $\overline{\text{Y}^1}$ in $\underline{P}$ x T

\margin{* Ask AG}

\clearpage
% EGA V, English translation by P. Blass, 1987 (cote 24): archive cote 24, site page 51 (PDF page 52).
% typescript. Draft compared: none.
% AI-transcribed by Claude Opus 5.5 (claude-opus-5-5, adaptive thinking, effort high), Message Batches API, batch ega24b (msgbatch_01RGJJeHdR96U16ZWFuUjtDx), 2026-11-13; not yet checked by eye.
\scanpage{24/51}
\mspage{51}

thus that the inverse image of $\underline{Y}^{1}$ (which is open in the above
closure) by h is an open subset of U, let us call it U1.  Because
of hg$_{1}$ = inclusion we see immediately that g$^{-1}$(U) $\to$ U is an open
immersion.  Since X is proper over k i.e. closed in $\underline{P}$, therefore g
is proper therefore surjective thus g$_{1}$ : g$^{-1}$(U) $\to$ U is
surjective, thus an isomorphism.  Also if W is an open subset of T
such that g$^{-1}$(W) $\to$ W is an isomorphism it follows that W is smooth
since $\underline{Y}^{1}$ is smooth, thus W \hand{$\subset$} U.  This proves 8.9.

The final assertions of 8.8 $\underline{Y}$sup sing = $\emptyset$ or dim $\underline{Y}$sup sing =
r - 2 and dim T = r - 1 are trivial:  the first one follows the
fact that $\underline{Y}$sing is irreducible of dim r and from the fact that Ysing
or Ysup sing [illegible]* is defined by the vanishing of a section D
of an invertible module.; the second from the fact that L being
finite over K we have deg tr L/k = deg tr K/k i.e. dim T = dim
$\underline{Y}$sing = r - 1.

\uline{Remark 8.10}

As we remarked in 8.9 with the notations of the corollary we
have g$^{-1}$(U) \hand{$\subset$} $\underline{Y}$sing - $\underline{Y}$sup sing but we notice that even if X is
closed in $\underline{P}$ this inclusion is not necessarily an equality, in
other words (noting that g$^{-1}$(U) is nothing else but the set of
points where g is étale, so that $\underline{Y}$sing - $\underline{Y}$sup sing is the set of
points where g is unramified) it may happen that there are points
y $\in$ $\underline{Y}$sing where g is unramified but not étale (which implies in
addition that g(y) is a point that is not geometrically normal and
even not geometrically unibranch of T).  In geometric terms this
corresponds to the following phenomenon; we may have a tangent

\margin{* ASK AG}

\clearpage
% EGA V, English translation by P. Blass, 1987 (cote 24): archive cote 24, site page 52 (PDF page 53).
% typescript. Draft compared: none.
% AI-transcribed by Claude Opus 5.5 (claude-opus-5-5, adaptive thinking, effort high), Message Batches API, batch ega24b (msgbatch_01RGJJeHdR96U16ZWFuUjtDx), 2026-11-13; not yet checked by eye.
\scanpage{24/52}
\mspage{52}

non-osculating hyperplane for X at a point x $\in$ X(k) such that
there exists another point x$^{1}$ $\in$ X(k) at which the same hyperplane
is tangent at x [or x$^{1}$ illegible]*. Indeed there are obvious
examples with $\underline{P}$ of dim two, X a non-singular curve of degree $\geq$ 4,
in any characteristic. [Note here: the ``double tangents'' of X
correspond to the double points of the ``dual curve.'']

\uline{Corollary 8.11}

Let us assume that k has characteristic zero. Then (a) The
image of $\underline{Y}$sup sing in $\underline{P}$v is of codimension $\geq$ 2. (b) the condition
(iii) of 8.8 is equivalent to other conditions i.e. the negation
of the other conditions, let us assume that $\underline{Y}$sing = Ysup sing], means
also that the image of $\underline{Y}$sing [or Ysup sing illegible]* in $\underline{P}$v is of
codimension $\geq$ 2.

Evidently, the assertion (b) implies (a) taking into account
8.8. But by dimension theory, (iii) means that L/K is a finite
extension (we could put it in the form ((iii) bis) in 8.8 and in
characteristic zero L is always separable over K hence the
condition (i bis) of 8.8.

\uline{Remark 8.12}

Geometrically the assertion (a) means essentially that for a
sufficiently general linear pencil of hyperplane sections every
member of the pencil is smooth or has for s ??? [illegible]* only
for geometric points singular points ordinary double points (and
in fact as one sees immediately it can be said in statement (a)
consequently in a form a little more precise--we have at most \uline{one}

\margin{* ASK AG (or figure out)}

\clearpage
% EGA V, English translation by P. Blass, 1987 (cote 24): archive cote 24, site page 53 (PDF page 54).
% typescript. Draft compared: none.
% AI-transcribed by Claude Opus 5.5 (claude-opus-5-5, adaptive thinking, effort high), Message Batches API, batch ega24b (msgbatch_01RGJJeHdR96U16ZWFuUjtDx), 2026-11-13; not yet checked by eye.
\scanpage{24/53}
\mspage{53}

such singular geometric point).  The assertion (b) means
essentially that if Ysing = $\underline{\text{Y}}$sup sing (which can be expressed
analytically by the vanishing of a certain section D of an
invertible module $W^{\otimes 2}{}_{X/k} \otimes O_{\underline{Y}}{}^{1}$, (1,1) over $\underline{\text{Y}}^{1}$), then for a
sufficiently general linear pencil of hyperplane sections all the
members of the pencil are smooth.  This second situation (whether
or not we are in characteristic zero) should entirely be
considered as exceptional.  [The variety Linl.... T = L illegible
handwriting on top of the page]*.  In the classical language it is
expressed, if there is no error, by saying that X is ruled for the
projective immersion considered [and if we so please] we have here
all that we need due to 8.5 and its corollaries to make explicit
and justify such a terminology in case if you feel inspired to
make connection with [la taupe Fr].  For example if dim X = 1 this
implies that X is a straight line [illegible]* x $\in$ X(k) so T
contains* ???.  (b) If the characteristic is p > \unsure{O|0}, we should
normally give examples (with dim $\underline{\text{P}}$ = 2, X a non-singular algebraic
curve) where the conditions of 8.6 are not satisfied i.e. $\underline{\text{Y}}$sing =
$\underline{\text{Y}}$sup sing and where nevertheless dim T = r - 1 i.e. examples where
L/K is a finite inseparable extension.  I am too lazy to construct
the examples but I do not doubt that such examples exist.

In (a) make a footnote to the following N° or paragraph where
we prove that if the exceptional 'ruled' case arises then by a
trivial modification of the projective immersion we find ourselves
again in the ``general'' situation of 8.8.

The part of the present section (or N$^{0}$) from 8.6 to here could
without a doubt be made into a separate section of a differential

\margin{* illegible}

\clearpage
% EGA V, English translation by P. Blass, 1987 (cote 24): archive cote 24, site page 54 (PDF page 55).
% typescript. Draft compared: none.
% AI-transcribed by Claude Opus 5.5 (claude-opus-5-5, adaptive thinking, effort high), Message Batches API, batch ega24b (msgbatch_01RGJJeHdR96U16ZWFuUjtDx), 2026-11-13; not yet checked by eye.
\scanpage{24/54}
\mspage{54}

character, whereas the beginning of our N$^0$ with the one that
follows should be merged together into a N$^0$ about the dimension of
exceptional [hyperplanes ???]*  I only use the fact that $\underline{Y}^{sing}$ has
dimension (r - 1)\hand{.} \uline{Proposition 8.13}\hand{.} We always assume that \margin{Display}
% NOTE: "Proposition 8.13" is circled by hand; a handwritten line runs from the circle to the circled marginal word "Display".
f : X $\to$ $\underline{P}$ is unramified and that X has no isolated points.

Let Z$_k$ be the part of $\underline{P}^\vee$ complement of the set of $\xi \in \underline{P}$ such that
$\phi_\xi$ should be $X_{k(\xi)}$ - regular and $\underline{Y}_\xi$ satisfies the geometric condition
R$_k$ then:

a) In order for Z$_{k-1}$ to be of codimension $\geq$ 2 in $\underline{P}^\vee$ it
suffices that every irreducible component X'$_i$ of X' of dimension $\leq$
K should be ruled for f.

b) In order to have Z$_k$ of codimension $\geq$ 2 in $\underline{P}^\vee$ it suffices
that every irreducible component X$_i$ of dimension $\leq$ K-1 (Ask A.G
illegible) should be made up of smooth points of X and should be
ruled.

Indeed for every $\xi$ geometrically singular the set of $\underline{Y}_\xi$
(NB:  We restrict ourselves to $\xi$ such that $\phi_\xi$ is $X_{k(\xi)}$ regular which
is harmless because of 8.2) and is the union of sing($\underline{Y}'_\xi$) and of
the inverse image T$_\xi$ of T in $\underline{Y}_\xi$ so that the codimension of this
singular set in $\underline{Y}_\xi$ is equal to the infimum of the codim (sing(Y'),
Y'$_\xi$) and of the codim (T$_\xi$, Y$_\xi$).  Let us restrict ourselves to $\xi$ such
that sing (Y'$_\xi$) is finite (which is harmless, this leads to place
ourselves in the complement of a set of codimension $\geq$ 2).  The
singular geometric points of $\underline{Y}'_\xi$ are therefore isolated.  The
conclusion follows easily from this and from 8.4.

Combining 8.13, 8.5 and 8.6 we find as usual.

\margin{* Ask AG}

\clearpage
% EGA V, English translation by P. Blass, 1987 (cote 24): archive cote 24, site page 55 (PDF page 56).
% typescript. Draft compared: none.
% AI-transcribed by Claude Opus 5.5 (claude-opus-5-5, adaptive thinking, effort high), Message Batches API, batch ega24b (msgbatch_01RGJJeHdR96U16ZWFuUjtDx), 2026-11-13; not yet checked by eye.
\scanpage{24/55}
\mspage{55}

\uline{Corollary 8.14}

We suppose that $f\colon X \to \underline{P}$ is unramified and that X has no
isolated points [illegible]* n.
a) suppose X is separable over k.  In order that the set of $\xi \in \underline{P}^{v}$
such that $\phi_{\xi}$ is $X_{k(\xi)}$ -regular and $\underline{Y}_{\xi}$ is separable, should have a
complement of codimension at least two it is necessary and
sufficient that every irreducible component $X_{i}$ of dimension one of
X should be formed from smooth points of X and should be ruled
relative to f and \uline{that} for every closed point x of X such that $\dim_{x}$
$X \geq 2$ we have prof x $X \geq 2$ $\mathrm{prof}_{x}X$, (conditions that are
automatically satisfied if X is geometrically normal and if all of
its irreducible components are of dim $\geq 2$).
b) let us assume that X is geometrically normal, in order that the
set of $\xi \in \underline{P}^{v}$ such that $\phi_{\xi}$ is $X_{k(\xi)}$ regular and $\underline{Y}_{\xi}$ is geometrically
normal should have a complement of codimension at least two it is
necessary and sufficient that every irreducible component $X_{i}$ of X
of dimension $\leq 2$ should be formed of smooth points of X and that
it should be ruled relative to f and that in addition for every
closed point x of X such that $\dim_{x} X \geq 3$ we have $\mathrm{prof}_{x} X \geq 3$.

\uline{Remark 8.15}

In 8.6, 8.13, 8.14 we make for X the hypothesis $(S_{k})$ (resp
$(R_{k})$) respectively : separable, respectively normal) that we wish
to recover as a conclusion for the hyperplane sections except
perhaps for $\xi$ from an exceptional set of codimension at least two.

This does not restrict the generality to tell the truth it
would have been better to get rid of this preliminary hypothesis)

\margin{* ASK AG}

\clearpage
% EGA V, English translation by P. Blass, 1987 (cote 24): archive cote 24, site page 56 (PDF page 57).
% typescript. Draft compared: none.
% AI-transcribed by Claude Opus 5.5 (claude-opus-5-5, adaptive thinking, effort high), Message Batches API, batch ega24b (msgbatch_01RGJJeHdR96U16ZWFuUjtDx), 2026-11-13; not yet checked by eye.
\scanpage{24/56}
\mspage{56}

since we see immediately with the help of results of par. 3.4 and
5.12 that if X does not satisfy the hypothesis in question, then
(by par. 5) if there exists a closed point x where the hypothesis
fails then for every $\xi$ such that $\phi_{\xi}$ is $X_{k(\xi)}$ regular condition that
only eliminate a set of codimention and such that $x \in Y_{\xi}$ (condition
that describes a set of exact codimension one), $\underline{Y}_{\xi}$ does not satisfy
the said hypothesis at x, the exceptional set $Z \langle \underline{P}^{v}$ is of
codimension one and not two.  (I may have somewhat exaggerated the
case $R_{k}$ where we still need some condition, $(S_{1})$ and perhaps of
equidimensionality perhaps.)
Remorse: In 8.13 and 8.14 it suffices to suppose that f: $X \rightarrow \underline{P}$ is
unramified at the smooth points of X; for the sufficiency part it
suffices only that they should be unramified over an open subset $X'$
of X where complement has codimension $\geq K + 1$.
\uline{Proposition 8.16}

Let us suppose f: $X \rightarrow \underline{P}$ unramified on an open subset
complementary to a set of codimension at least two, X
geometrically normal and of depth at least three at its closed
points, finally X geometrically integral and proper over k.  Then
the set of $\xi \in \underline{P}^{v}$ such that $\underline{Y}_{\xi}$ is geometrically normal and
geometrically integral of dimension equal to dim X-1 (is
constructible and) has a complement of codimension at least two.

Indeed by 8.14 b) such is the case for the property ``$\underline{Y}_{\xi}$ is
geometrically normal of dimension dim X-1'' (the dimensional
property expresses that $\phi_{\xi}$ is $X_{k(\xi)}$ - regular).

\clearpage
% EGA V, English translation by P. Blass, 1987 (cote 24): archive cote 24, site page 57 (PDF page 58).
% typescript. Draft compared: none.
% AI-transcribed by Claude Opus 5.5 (claude-opus-5-5, adaptive thinking, effort high), Message Batches API, batch ega24b (msgbatch_01RGJJeHdR96U16ZWFuUjtDx), 2026-11-13; not yet checked by eye.
\scanpage{24/57}
\mspage{57}

Therefore, under 6.1 all the $Y_\xi$ are geometrically connected.
Since $\underline{Y}_\xi$ is geometrically normal it is geometrically integral if
and only if it is geometrically connected which gives the proof.

\uline{Remarks: 8.17}

a) The hypothesis of 8.16 implies that dim $X \geq 3$. and that
dim $f(x) \geq 3$ (without the hypothesis of normality and of non-%
ramification) the set of $\xi$ such that $\underline{Y}_\xi$ is geometrically
irreducible has a complement of codimension at least two. We can
prove in every case that it does not contain a hyperplane of below
b) the conclusion of 8.16 is false if we leave out the assumption
that $\mathrm{prof}_x X \geq 3$ for $x$ closed, for example it is false for a non-%
singular quadric $X$ in $\underline{P}^3$ [illegible]*
the hyperplane sections are reducible (in fact formed by pairs of
concurrent lines) and they form therefore a two dimensional family
thus of codimension one in $\underline{P}$ (indeed the dual of the quadric is a
quadric in the dual space relative to the dual form...)

In the case of a non-singular surface in a projective space
this situation however should be considered exceptional of the
following section (or $N^0$). Let us suppose [illegible]* that $X$ is
smooth of dimension two [illegible]* integral proper over $k$ and $f$
an immersion. Then it follows from 6.1 and 8.8 and 8.14 that if
$Y^{\mathrm{sing}} \to \underline{P}^\vee$ is not generically finite and inseparable one of the
(???) $\xi \in \underline{P}^\vee$ such that $\underline{Y}_\xi$ is separable over $k(\xi)$ with at most two
irreducible [illegible]* a complement of codim at least two.

\margin{* Ask AG}

\clearpage
% EGA V, English translation by P. Blass, 1987 (cote 24): archive cote 24, site page 58 (PDF page 59).
% typescript. Draft compared: none.
% AI-transcribed by Claude Opus 5.5 (claude-opus-5-5, adaptive thinking, effort high), Message Batches API, batch ega24b (msgbatch_01RGJJeHdR96U16ZWFuUjtDx), 2026-11-13; not yet checked by eye.
\scanpage{24/58}
\mspage{58}

We shall now examine more precisely the case of surfaces (the
case of curves does not arise evidently, from the point of view of
irreducibility of hyperplane sections).
NB  I noticed with fright that the quadric is not entitled to the
use [Fr il l'emploi que j'ai fait de mot
``reglé,'' à être appelée réglée!]*
This is in disagreement with our fathers and it would be necessary
to invent a more adequate word for the notion used here.

\uline{Proposition 8.18}

Let us suppose that k is algebraically closed, X is integral
(respectively integral and normal) of dimension $\geq 2$ and proper
over k, let T be a closed finite subset of X such that X - T is
smooth and let f/X - T be unramified.  In order that the set of $\xi$
$\in P^{\vee}$ such that $Y_{\xi}$ should be geometrically irreducible (respectively
geometrically integral) of dimension 1 should have a complement of
codimension $\geq 2$ it is necessary and sufficient that the following
two conditions should be satisfied:  a) For every $x \in T$ there
exists a hyperplane section $Y_{\xi}$ ($\xi \in P(k)$ ) passing through x of
dimension d - 1 and which is irreducible, b) X' = X - T is ``ruled''
(\uline{sic}) for f \uline{or} there exists a hyperplane section $Y'_{\xi}$ ($\xi \in P^{\vee}(k)$ ) of
X' which is of dimension d - 1 (\uline{non}??? singular and irreducible.
[illegible ask AG]

Let us first assume that X is geometrically normal.  We have
already seen then by 8.14  a) that we can find a closed \uline{subset} $Z^{1}$
of P of codim $\geq 2$ such that $\xi \in P-Z^{1}$ implies that $Y_{\xi}$ is separable
over $k(\xi)$ and of dimension d-1 for such a $\xi$, it amounts to the same

\margin{* ASK AG}

\clearpage
% EGA V, English translation by P. Blass, 1987 (cote 24): archive cote 24, site page 59 (PDF page 60).
% typescript. Draft compared: none.
% AI-transcribed by Claude Opus 5.5 (claude-opus-5-5, adaptive thinking, effort high), Message Batches API, batch ega24b (msgbatch_01RGJJeHdR96U16ZWFuUjtDx), 2026-11-13; not yet checked by eye.
\scanpage{24/59}
\mspage{59}

that $Y_{\xi}$ should be geometrically irreducible or geometrically
integral, and the two problems [(\uline{respé et non respé}) Fr p 50]
considered in 8.18 are therefore equivalent. On the other hand,
by 5.6, the set U of $\xi \in P$ such that $Y_{\xi}$ is geometrically integral
of dim d-1 (the dimension hypothesis stating that $\phi_{\xi}$ is $X_{k(\xi)}$
regular) is \uline{open}. We will exhibit a non-empty open \uline{evident} subset
P-Z contained in U and taking for Z the \uline{union of} $g(\underline{Y}^{1\ \mathrm{sing}}$ and of
the hyperplanes $H_{x}$ of $P^{v}$ defined by the f(x), $X \in T$. For $\xi \in P^{v}$ -
Z, $Y_{\xi}$ is smooth of dimension
(d - 1) and since it is geometrically connected by 6.1 it is
geometrically integral. We have to, therefore, express (explain)
(prove) that every irreducible component of codimension one of Z
meets the open set U. But these irreducible components are the $H_{x}$
[they are repeated possibly, but it is not essential] and also $g(Y^{1}$
$^{\mathrm{sing}})$ when the latter are indeed of codimension one i.e. $X^{1}$ ``not
ruled'' for f (Nota Bene: we use the irreducibility of $Y^{1\mathrm{sing}}$). On
the other hand, in order that this latter set should meet the open
set U it is necessary and sufficient that $g(Y^{1\mathrm{sing}})$ which contains
an open and dense set) should meet U. This proves 8.18 in this
case. If we do not suppose that X is normal, we apply the
previous result to the normalization of X the reasoning is
immediate and I do not give the details here. N.B. In the \uline{case}
[respé] 8.18 is contained in 8.16 more precisely except in the
case d = 2. It is for the case \uline{[non-respé] that it may be better}
\uline{not to require]} d = 2 and not only $d \geq 2$ .....

It remains to explain, (make explicit) the conditions a) and
b) of 8.18. This leads us to exmaine in a general way the

\clearpage
% EGA V, English translation by P. Blass, 1987 (cote 24): archive cote 24, site page 60 (PDF page 61).
% typescript. Draft compared: none.
% AI-transcribed by Claude Opus 5.5 (claude-opus-5-5, adaptive thinking, effort high), Message Batches API, batch ega24b (msgbatch_01RGJJeHdR96U16ZWFuUjtDx), 2026-11-13; not yet checked by eye.
\scanpage{24/60}
\mspage{60}

following situation.  We suppose that $X$ is geometrically
irreducible over $k$ and we (give ourselves) consider a linear
subvariety $L$ of $P$ (corresponding to the question of studying the
hyperplane sections of $X$, passing through a given point $x$ or
tangent to $X$ at a given smooth point), formed therefore, by the
hyperplane containing a linear subvariety $L$ of $P$ (resp. a point,
or the image of a tangent space to $X$ at a smooth point in the two
cases considered) and we ask the question [de ravoi] if for the
generic point of $L$ (therefore for all the points of a non-empty
open subset of $L$) $Y$ is geometrically irreducible of dim = dim $X$ -
1.  This is a variant of Bertini's theorem, which [(devrait
figurer)51] must appear in No. 3, and is treated by exactly the
same method, [(ou, si on veut, s'y ramène) Fr 51].  The dimension
question is simply stated for $f(x)$ C $L^0$ i.e. if $X^1 = f^{-1}(P - L^0)$ is
a dense open subset of $X$.  Let $Q$ be the projective space of
hyperplanes passing through $L^0$ (N.B. if $L^0$ is defined by a vector
subspace $F^0$ of $E$ we have $Q = P(F^0)$ and we consider the canonical
morphism (deduced from $F^0 \to E$, cf.  Chap. II).
\[ U : \underline{P} - L^0 \to Q \]
and we consider
\[ g = Uf^1 : f^{-1}(P - L) = X^1 \to Q \]
so that $L \simeq \check{Q}$ and the family of $X^1{}_{\hand{$\xi$}}$ ($\hand{$\xi$} \in L$) is nothing else than
the family of hyperplane sections relative to the morphism $g$.  On
the other hand, we see immediately that for every $\hand{$\xi$}\in L$, ``\uline{general}''
$X^1$ is dense in $X$, so that $X^1$ is geometrically irreducible if and

\clearpage
% EGA V, English translation by P. Blass, 1987 (cote 24): archive cote 24, site page 61 (PDF page 62).
% typescript. Draft compared: none.
% AI-transcribed by Claude Opus 5.5 (claude-opus-5-5, adaptive thinking, effort high), Message Batches API, batch ega24b (msgbatch_01RGJJeHdR96U16ZWFuUjtDx), 2026-11-13; not yet checked by eye.
\scanpage{24/61}
\mspage{61}

only if X is such.  This assumed, the theorem of Bertini-Zariski
shows us that we have the wanted conclusion of irreducibility
provided that dim $g(X^1) \geq 2$.  (To tell the truth one could give a
reciprocal to 3.1 as follows:  If X is geometrically irreducible Y
is geometrically irreducible if and only if either dim $f(X) = 2$ \uline{or}
dim $f(X) = 1$ and $f(X)$ is contained in a straight line D defined
over k and the generic fiber of $X \to D$ is geometrically
irreducible).  This also allows us in the present version with L
to have a necessary and sufficient condition of geometric
irreducibility of $Y_{\xi}$, $\xi$ generic in L.

From the [(\uline{cunutesque}?) Fr] point of view and in terms of
field theory we can express the condition in terms of
transcendence degree in the following fashion.  We choose a
``hyperplane at infinity'' containing neither $L^{\circ}$ nor X and we place
ourselves in its complement i.e. over a scheme of affine type
essentially.  We choose a basis of the space of linear forms
vanishing on L, let it be $T_1, \ldots, T_p$ (p = codim ($L^{\circ}$, P) and we
consider their inverse images $t_1,, \ldots., t_p$ in the the field of
fractions K of X (  X assumed integral).  At least one of the $t_i$,
let us say $t_1$ is $\neq 0$.  Let us consider therefore, $a_1 = t_2/t_1, \ldots.,$
$a_{p-1}) = tp/t_1$ then dim $g(X^1)$ is nothing else but the transcendence
degree of $K(a_1, \ldots\ a_{\,p-1})$ \hand{$\subset$} K over k.  Therefore if the
transcendence degree is $\geq 2$ we are o.k.  If it is one then we must
require that over k, $f(X)$ is contained in a linear subvariety of P
containing $L^{\circ}$ and of dimension at most one and that the \uline{generic}
fiber of $g{:}X^1 \to g(X^1)$ should be geometrically irreducible.

\clearpage
% EGA V, English translation by P. Blass, 1987 (cote 24): archive cote 24, site page 62 (PDF page 63).
% typescript. Draft compared: none.
% AI-transcribed by Claude Opus 5.5 (claude-opus-5-5, adaptive thinking, effort high), Message Batches API, batch ega24b (msgbatch_01RGJJeHdR96U16ZWFuUjtDx), 2026-11-13; not yet checked by eye.
\scanpage{24/62}
\mspage{62}

Let us suppose that $L^{\circ}$ is of dimension q, so that the fibers
of $u\colon P - L^{\circ} \to Q$ are of dimension $q + 1$ so that those of g are dim
$\leq q + 1$, and consequently we have $\dim g(X^{1}) \geq \dim f(X) - (q + 1)$ so
that the dimension condition for $g(X^{1})$ is verified in view of the
fact that $\dim f(X) \geq q + 3$. If $q = 0$ we find the fact indicated
in 8.17 a). Returning to conditions of 8.18 we see that condition
a) relative to an $x \in T$ is satisfied provided x is not ``conical at
x relative to f'' in an obvious sense. Maybe it will be better to
introduce these latest \uline{Bertinisque} developments in the next
section... change of projective immersion.

\clearpage
% EGA V, English translation by P. Blass, 1987 (cote 24): archive cote 24, site page 63 (PDF page 64).
% typescript. Draft compared: none.
% AI-transcribed by Claude Opus 5.5 (claude-opus-5-5, adaptive thinking, effort high), Message Batches API, batch ega24b (msgbatch_01RGJJeHdR96U16ZWFuUjtDx), 2026-11-13; not yet checked by eye.
\scanpage{24/63}
\mspage{63}

\uline{Section 9}

\uline{Change of Projective Embedding}

\uline{9.1}

For every integer $n > 0$ let $P(n) = P(\mathrm{Sym}^{n}\ (E))$, we have an
evident immersion $u_{n}\colon P \to P(n)$, since $O\ P(n)$ is generated by its
sections over every open affine of $S$ and that $P_{*}(O_{p}(n))\ \underline{\mathrm{Sym}}^{n}(\underline{E})^{*}$
(???) where $p\colon P \to S$ is the projection.  If $f\colon X \to P$ is an
unramified morphism (resp. an immersion) it is the same with $u_{n}f\colon X$
$\to P(n)$.  There is sometimes an advantage in the study of $X$ in
replacing $f$ by $u_{n}f$ in order to avoid a very special behavior and
sometimes embarrassing f \ \ a certains regards.  (An example of
such peculiarity is the one indicated (\uline{sic}) in 8.12 b), where $\underline{Y}^{\mathrm{sing}}$
$\to P^{v}$ has an image of dimension $r - 1$ but gives rise to an
inseparable extension of fields.  Another one is that given by the
quadric surfaces in $P^{3}$ to know that all the singular hyperplane
sections are geometrically reducible.  (in spite of the fact that
$X$ is geometrically irreducible.)

\uline{Proposition 9.2}

We suppose $S = \mathrm{Spec}(k)$, $X$ smooth over $k$, and $f\colon X \to P$
unramified.  Let $n \geq 2$ and let us consider $f_{n} = u_{n}f$.  Then $f_{n}\colon X \to$
$P(n)$ satisfies the equivalent conditions of 8.8, in particular for
$\xi \in P(n)$ in the complement of a set of codimension $\geq 2$, the
corresponding hyperplane section $Y_{\xi}$ is smooth \uline{or (admits)} only a
finite number of non-smooth points which are geometrically

\margin{* illegible}

\clearpage
% EGA V, English translation by P. Blass, 1987 (cote 24): archive cote 24, site page 64 (PDF page 65).
% typescript. Draft compared: none.
% AI-transcribed by Claude Opus 5.5 (claude-opus-5-5, adaptive thinking, effort high), Message Batches API, batch ega24b (msgbatch_01RGJJeHdR96U16ZWFuUjtDx), 2026-11-13; not yet checked by eye.
\scanpage{24/64}
\mspage{64}

ordinary singularities.  If $f$ is an immersion there is at most one
such singular point and it is rational over $k(\xi)$.

N.B. One would have to announce 8.8 in a manner such as not
to exclude the case where $f$ is not an immersion.  The verification
is essentially trivial under the condition (ii bis) of 8.8.
Without a doubt we should make explicit in 9.1 that the hyperplane
sections of $X$ relative to $f_n$ are nothing else the sections of $X$
relative to $f_n$ are nothing else but the ``sections'' of $X$ by
hypersurfaces of degree $n$ in place of hyperplanes.

\uline{Proposition 9.3}

Suppose that $X$ is geometrically irreducible and $\dim f(X) \geq 2$,
let $x \in X(k)$ let $n \geq 2$ .

a) Let us consider the linear family of hyperplane sections
of $X$ relative to $f_n =u_nf$ \uline{which pass through} $x$, its generic element
defines a $Y_\xi{}^{(n)}$ which is geometrically irreducible.

b) Let $L$ be a linear subvariety of $P$ passing through $f(x)$ and
not containing $f(X)$.  Let us consider the linear family of
hyperplane sections of $X$ relative to $f_n$ defined by the hyperplanes
of $P(n)$ ``tangent to $L$ at $x$'' (i.e. defined by the $n$-forms over $P$
which over $L$ are zero of order at least two at $x$) , its generic
member is a $Y_\xi{}^{(n)}$ which is geometrically irreducible.

c) Let us suppose that $X$ is smooth at $x$ and that $n \geq 3$ where
$f(X)$ is not contained in a plane defined over $k$.  Let us consider
the family of hyperplane sections $Y^{(n)}$ of $X$ relative to $f_n$ which
are ``tangent to $X$ at $x$''.  Then the generic member of the latter
defines a $Y_\xi{}^{n)}$ that is geometrically irreducible.

\clearpage
% EGA V, English translation by P. Blass, 1987 (cote 24): archive cote 24, site page 65 (PDF page 66).
% typescript. Draft compared: none.
% AI-transcribed by Claude Opus 5.5 (claude-opus-5-5, adaptive thinking, effort high), Message Batches API, batch ega24b (msgbatch_01RGJJeHdR96U16ZWFuUjtDx), 2026-11-13; not yet checked by eye.
\scanpage{24/65}
\mspage{65}

The proof is essentially trivial in terms of the criteria of
the end of the previous section. Taking an affine model of $P$
containing $f(X)$ we are reduced a) to finding 3 polynomials in the
coordinates $T_1$, $T_2$,....., $T_s$ of degree $\leq 2$, let them be $P$, $Q$, and $R$
such that $Q(t)/P(t)$ and $R(t)/P(t)$ are algebraically independent
over $k$ in $K$ (where $K$ is the function field of $X$ and $t = (t_1$,....,
$t_r)$ is the system of elements of $K$ defined by the $T_i$); in b) we
require also that $P$, $Q$, and $R$ should vanish to order two at least
on $L$ which we can in addition suppose to be defined by the
equations $T_1$,...., $T_s = 0$; finally in c) it is the same but $L$ is
the image of the tangent space to $X$ at $x$ and we allow possibly to
take $P$, $Q$, and $R$ of degree 3 i.e. move margin, the hypothesis that
dim $f(X) \geq 2$ signifies that the transcendence degree of $K(t_1$,
$t_2$,...., $t_r)$ over $k$ is $\geq 2$ i.e. we can find $t_1$, $t_2$ let us say
algebraically independent. In a) we take therefore $P = T_1$, $Q =$
$T_1{}^2$, $R = T_1T_2$, in b) we (analogously) do the same noting that, we
may there choose $t_1$ leading to $T_1$ zero over $L$ due to the fact that
$f(X)$ C $L$ (which implies that there exists an index $i$ between one
and $s$ such that $t_s \neq 0$, so that $t_s$ is not a constant (since $t_s$ is
zero at $x$) therefore $t_s$ is not algebraic over $k$ (\uline{???}) (N.B. we may
suppose $k$ algebraically closed). The case c) follows from b)
except in the case where $f(X)$ is contained in the image of $L$ by $f$
of the tangent space to $X$ at $x$\ \ \ . [If dim $X = 2$ this case is
effectively exceptional (the trace of a quadric surface tangent to
a plane on that plane is in general formed by two intersecting
lines and is therefore not irreducible]. But to treat that case,
in the forms $P$, $Q$, and $R$ made explicit above we may replace

\clearpage
% EGA V, English translation by P. Blass, 1987 (cote 24): archive cote 24, site page 66 (PDF page 67).
% typescript. Draft compared: none.
% AI-transcribed by Claude Opus 5.5 (claude-opus-5-5, adaptive thinking, effort high), Message Batches API, batch ega24b (msgbatch_01RGJJeHdR96U16ZWFuUjtDx), 2026-11-13; not yet checked by eye.
\scanpage{24/66}
\mspage{66}

evidently X by L itself, where the solution is trivial.  (If dim L
= 2 take P = $T_r{}^2$, Q = $T_r{}^3$, R = $T_r{}^2T_{r-1}$ and note that Q/R = $T_r$ and R/P
= $T_{r-1}$ are linear forms independent over L.  Therefore
algebraically independent.  If dim L $\geq$ 3 then $T_{r-2}$, $T_{r-1}$, $T_r$
linearly independent over L and we take
\[ \underline{P} = {T^2}_r, \ Q = {T^2}_{r-1}, \ R = {T^2}_{r-2} \]
Conjugating with 8.18 we find a

\uline{Corollary 9.4}

Finally we must combine the latter with 9.2 in order to find
a recapitulating theorem in the ``excellent case.''

\uline{Theorem 9.5}  (illegible page 52) or 55.*

(If X is smooth and proper and geometrically integral over k and
f:X $\rightarrow$ P is unramified X of dimension $\geq$ 2.  (Ask Grothendieck).

\uline{Section 10}

\uline{10.0  Pencils of hyperplane sections and fibrations of blown up}
\uline{varieties.10.1}

Let Z be the P-exceptional set in $P^{v}$ relative to a
constructible property P such that Z is a constructible part of $P^{v}$.
Let us suppose S=Spec(k).  We see (cf. No. 12 where we catch up
with things which should have come in without a doubt in (previous
No.'s) that in order to have codim (Z, $P^{v}$) $\geq$ 2 it is necessary and
sufficient that ``every sufficiently general line'' L in P should

\margin{* Ask AG}

\clearpage
% EGA V, English translation by P. Blass, 1987 (cote 24): archive cote 24, site page 67 (PDF page 68).
% typescript. Draft compared: none.
% AI-transcribed by Claude Opus 5.5 (claude-opus-5-5, adaptive thinking, effort high), Message Batches API, batch ega24b (msgbatch_01RGJJeHdR96U16ZWFuUjtDx), 2026-11-13; not yet checked by eye.
\scanpage{24/67}
\mspage{67}

not meet Z or also again Z and it suffices that there should exist
one (a single) L in P not meeting k (\uline{Z AG's error P.B.})*. If k is
infinite it is necessary and sufficient that there should exist a
single straight line L in P that does not meet Z. We call a
``linear pencil of hyperplane sections of x'' defined by the
straight line L in P the L prescheme $Y_L$ (definition valid for any
S). Then the previous reflections together with results of No's.
8 and 9 give us criteria for the existence of such pencils having
the fibers $\underline{Y}_\xi$ ($\xi \in L$) all satisfying the property P first of all in
the case where S is an infinite base field. Taking into account
8.2, if for every associated prime cycle on X we have dim f(T) $> 0$
we can (by taking the property $P^1 = P$ + condition of regularity for
$\phi_\xi$) require that the pencil $Y_L$ should be flat over L. In the case
where S is arbitrary we can again often proceeding by the
procedure of 7.1 construct such a pencil over an open neighborhood
of a given point s of S in view of the fact that k(s) is infinite
and we should know that Z is closed (which is assured in diverse
cases by the results of Par. 5 and the assumption that $X \to S$ is
proper). To do it right it would be convenient after general
explanation of this type to give recapitulating statements where
we effectively apply the preceding results for a certain number of
properties of this nature (and also comprising module properties).
As a minimum in this sense we must give here the reformulation of
9.5 in terms of linear pencils - a fact constantly used in
geometric applications.

\margin{* Ask A.G.}

\clearpage
% EGA V, English translation by P. Blass, 1987 (cote 24): archive cote 24, site page 68 (PDF page 69).
% typescript. Draft compared: none.
% AI-transcribed by Claude Opus 5.5 (claude-opus-5-5, adaptive thinking, effort high), Message Batches API, batch ega24b (msgbatch_01RGJJeHdR96U16ZWFuUjtDx), 2026-11-13; not yet checked by eye.
\scanpage{24/68}
\mspage{68}

\uline{10.2}

By a polarity, to a straight line L in $P^{\vee}$ there corresponds a
linear subvariety $L^{\circ}$ of codim 2 in P (S arbitrary).  Let us put T =
$X \times_{P} L^{\circ}$.  Another way to describe T is as follows; L is defined by
a locally free quotient of rank 2 of $\underline{E}^{\vee}$ or what is the same by a
submodule, locally a direct factor, $\underline{F}$ of $\underline{E}$ everywhere of rank two.
Let us consider the composed homomorphism
\[ F_{X} \to E_{X} \to O_{X}\,(1), \]
then T is nothing else but the scheme of zeroes of this composed
homomorphism or what is the same it is defined by the ideal J,
image of the corresponding homomorphism (obtained by twisting by
$O_{X}(-1)$)
\[ : F_{X}(-1) \to O_{X} \;. \]

Let us suppose that this homomorphism is regular which means
that if we write down locally a totally ordered basis of $F_{X}(-1)$
then its image in $O_{X}$ forms an $O_{X}$ - regular sequence, a condition
that does not depend on the basis chosen and that can be announced
intrinsically also by saying that $F_{X}(-1) \otimes O_{X}/J \to J/J^{2}$ is an
isomorphism and that $V(J) = T \to X$ is a regular immersion [NB: we
should somewhere pull out of find out the general situation] with
a homomorphism $G \to O_{X}$, G locally free over the prescheme X, for
example in the section about regular immersions] we have then

\clearpage
% EGA V, English translation by P. Blass, 1987 (cote 24): archive cote 24, site page 69 (PDF page 70).
% typescript. Draft compared: none.
% AI-transcribed by Claude Opus 5.5 (claude-opus-5-5, adaptive thinking, effort high), Message Batches API, batch ega24b (msgbatch_01RGJJeHdR96U16ZWFuUjtDx), 2026-11-13; not yet checked by eye.
\scanpage{24/69}
\mspage{69}

\uline{Theorem 10.2}

With the above hypothesis the linear pencil $Y_L$ with the
canonical projection $Y_L \rightarrow X$ is X isomorphic in a unique fashion to
the blow-up of prescheme X with center T.

To understand the meaning of this theorem it is convenient to
notice at the beginning of the section that if S = Spec(k) then
for a 'sufficiently general' straight line L in P the condition of
regularity is verified (cf. \uline{catching up indicated in No. 1}) namely
by 5.3.  In what follows in the construction of good linear
pencils.  (\uline{above five lines a mystery}).*
indicated or anticipated at the beginning of the present N°, namely
[illegible]* 5.3.  In what follows in the construction of the
``good'' linear pencils indicated in the beginning of the N°, we
could require that the described pencil should satisfy the said
condition (which is a condition of the same type but different
from the one that consists in requiring that for every $\xi \in L$, $\phi_\xi$
should be $X_{k(\xi)}$ - regular).  We should include the condition in
question in the proposed recapitulating statements.

On the other hand, practically 10.2 is used only in the
situation of 9.5, which makes desirable to announce the
reformulation of 9.5 in terms of pencils which according to 10.2,
in order to be able to include in the statement in question also
the isomorphism of the pencil with a blow up (i.e. to give a
description of the situation permitting a suitable reference).  We
obtain thus a way for every projective smooth geometrically
connected X of dim $\geq 2$ over an infinite field k, to find a closed
smooth subscheme non-empty of codimension two at every one of its

\margin{* Ask AG}

\clearpage
% EGA V, English translation by P. Blass, 1987 (cote 24): archive cote 24, site page 70 (PDF page 71).
% typescript. Draft compared: none.
% AI-transcribed by Claude Opus 5.5 (claude-opus-5-5, adaptive thinking, effort high), Message Batches API, batch ega24b (msgbatch_01RGJJeHdR96U16ZWFuUjtDx), 2026-11-13; not yet checked by eye.
\scanpage{24/70}
\mspage{70}

points such that the blown up scheme admits a fibration over $\underline{P}^1$,
with all the fibers are geometrically integral and such that all
the fibers are smooth except at most a finite number, the latter
having at most a geometrically singular point and such a point
being rational over k and geometrically an ordinary singularity.

This explains the importance of a deep study just [58]
started at the present time) of such fibrations with singular
fibers, to reduce in a certain measure (to some extent) the study
of projective smooth varieties of dimension d to those of
(families depending on one parameter) of projective varieties of
dimension (d - 1) that may have ordinary singularities.

The statement 10.2 is a more or less immediate consequence of
the following which is completely independent of the story of
hyperplane sections and would be without a doubt better in its
place in a ${}^{\mathrm{V}}$N$^{\circ}$ added ``regular immersions.''*
Proposition 10.3 is crossed out.  Ask AG if that is his
intention.*

\uline{Corollary 10.4} (of 10.2)

Let X be a smooth curve geometrically connected in a
projective space P over an infinite field k and let n $\geq$ 2.  Then
there exists a linear pencil of n- forms over P defining a
morphism X $\to$ P$^1$ having the following property:  the morphism is
generically étale of degree d and for every geometric point s of
P$^1$, X$_s$ is étale over the algebraically closed field k$'$ = k(s) \uline{or} it
is k$'$ isomorphic to the sum of d - 2 schemes Spec k$'$ and the scheme

\margin{* Ask Grothendieck}

\clearpage
% EGA V, English translation by P. Blass, 1987 (cote 24): archive cote 24, site page 71 (PDF page 72).
% typescript. Draft compared: none.
% AI-transcribed by Claude Opus 5.5 (claude-opus-5-5, adaptive thinking, effort high), Message Batches API, batch ega24b (msgbatch_01RGJJeHdR96U16ZWFuUjtDx), 2026-11-13; not yet checked by eye.
\scanpage{24/71}
\mspage{71}

$I_k{}' = \mathrm{Spec}\ k'[t]/_{(t^2)}$.  In the language of the forefathers:  there
is at most one point of ramification and it is ``quadratic.''

\uline{Section 11}

\uline{Grassmanians}

Since we will now use linear subvarieties of P not only of
relative dimension 0 and n - 1 it is clear that we shall need some
notions about grassmanians and some ['\uline{sorites}']** (facts) of the
nature of 'elementary geometry' about the constructions concerning
linear varieties which should all come at the begining of the
paragraph.  In addition \uline{one takes in practice sometimes} any linear
sections and not only hyperplane sections and it is proper to
review (revisit) in this enlarged spirit all the previous No's.

Let E be a quasi-coherent module over the prescheme (\uline{X or S})
and let n be an integer > 0.  Let us consider the functor $(\mathrm{Sch})^{\circ}/_S$
$\rightarrow$ (Ens) defined by $\underline{\mathrm{Grass}}_n\ (E)\ (S^1)$ = set of quotient modules,
locally free and of rank n of $E_S{}^1$.

This functor is representable and the prescheme over S which
represents it will also be denoted $\underline{\mathrm{Grass}}_n\ (E)$.  To prove the
representability consider the natural homorphism of functors
\[ \underline{\mathrm{Grass}}_n\ (E) \rightarrow \underline{\mathrm{Grass}}_1\ (nE) = P(nE) \]
defined by associating with every locally free quotient of rank n,
G of $E_S{}^1$ the locally free module of rank one G considered as a
quotient of $E_S{}'$.  We prove as in \uline{Seminaire Cartan} that this
morphism is ``representable by a closed immersion'' (\uline{for closed}
\uline{immersions}) such that $\underline{\mathrm{Grass}}_n\ (E)$ appears as a closed subscheme of
P(nE); in particular it is separated over S and quasi-compact over

\margin{** what is best translation of this word?}

\clearpage
% EGA V, English translation by P. Blass, 1987 (cote 24): archive cote 24, site page 72 (PDF page 73).
% typescript. Draft compared: none.
% AI-transcribed by Claude Opus 5.5 (claude-opus-5-5, adaptive thinking, effort high), Message Batches API, batch ega24b (msgbatch_01RGJJeHdR96U16ZWFuUjtDx), 2026-11-13; not yet checked by eye.
\scanpage{24/72}
\mspage{72}

S and if E is of finite type it is projective over S.  If E is of
finite presentation then that is also the case for $\underline{Grass}_n(E)$.
Indeed we may suppose that S is affine S = Spec(A) so that E comes
from a module of finite type over a subring of finite type of A -
since the formation of $\underline{Grass}_n(E)$ is evidently compatible with base
change over S.

Since E is locally free therefore $\underline{Grass}_n(E)$ is smooth over S
with geometrically connected fibers.  This comes from a more
precise fact:  If E is free of rank r then $\underline{Grass}_n(E)$ may be
recovered from $\binom{\hand{$r$}}{n}$ open subsets each one of which is S isomorphic
to affine space of relative dimension n(r-n) over S.  This
decomposition corresponds to the choice, thanks to the base of E
to $\binom{\hand{$r$}}{n}$ decompositions of E by exact sequences (s) $0 \to E' \to E \to$
$E'' \to 0$ with E' locally free of rank n.  Such an exact sequence
allows us to define a sub-functor.  $\underline{Grass}_n(s)$ of $Grass_n(\underline{E})$ by
limiting ourselves to quotients $\underline{G}$ of $E_S{}'$ locally free of rank n,
such that the composed homomorphism $\underline{E}'_S{}' \to \underline{E}_S{}' \to \underline{G}$ should be
surjective (therefore bijective).  But the inclusion $\underline{Grass}_n(S) \to$
$\underline{Grass}_n(E)$ is representable by open immersion and on the other hand
$\underline{Grass}_n(E)$ is representable by open immersion and on the other hand
$\underline{Grass}_n(S)$ is canonically isomorphic to the fiber bundle $V(\underline{Hom}_{O_S}$
$(\underline{E}', \underline{E}''))$.

As a result for example of this particular structure we may
mention that if $s \in S$, then ($\underline{E}$ being locally free of finite rank)
every point of $\underline{Grass}_n(\underline{E})$ with values in k(s) lifts to a section
over a neighborhood of s.  On the other hand if S = Spec(k), k an
infinite field, then every open non-empty subset of $\underline{Grass}_n(E)$

\clearpage
% EGA V, English translation by P. Blass, 1987 (cote 24): archive cote 24, site page 73 (PDF page 74).
% typescript. Draft compared: none.
% AI-transcribed by Claude Opus 5.5 (claude-opus-5-5, adaptive thinking, effort high), Message Batches API, batch ega24b (msgbatch_01RGJJeHdR96U16ZWFuUjtDx), 2026-11-13; not yet checked by eye.
\scanpage{24/73}
\mspage{73}

contains a k-rational point.  A point of $\underline{\mathrm{Grass}}_{n}(E)$ with values in S
i.e. a locally free quotient module G of rank n of E canonically
defines a subscheme of P(E), (i.e., to say) P(G).  Such a
subscheme (but without \uline{imposing} the rank of G) is called a \uline{linear}
\uline{subvariety} of P(E) (relative to S if there is a possibility of
confusion).  It is therefore a projective fibration of relative
dimension (n-1) if $n \geq 1$, (and empty is $n = 0$).  We immediately
verify that the section of $\underline{\mathrm{Grass}}_{n}(E)$ i.e. G is known if we know the
linear subvariety corresponding to P(E).  In this manner the
grassmanian can be interpreted as representing the functor
(``linear subvarieties of relative dimension n - 1 of $\underline{P}_{S}{}^{1}$ '') for $S^{1}$
variable $n \geq 1$.  It is furthermore possible to give an intrinsic
characterization of the latter functor i.e. of the notion of
linear subvariety of relative dimension m and that are of
``projective degree one'' at every $s \in S$; this characterization will
be given in a later chapter and we shall not need it at all here.

Let us again suppose that E is locally free of rank r, let $E^{\vee}$
be its dual.  Then by a polarity we find a canonical isomophism
$\underline{\mathrm{Grass}}_{n}(E) \simeq \underline{\mathrm{Grass}}_{r-n}(E^{\vee})$ that assigns to a quotient $\underline{G}$ of $\underline{E}$ the
quotient $\underline{E}^{\vee}/\underline{G}^{\vee}$ of $\underline{E}^{\vee}$.  From the point of view of linear varieties
to a linear variety L of relative dimension m of P there
corresponds the linear dual variety $L^{\circ}$ of relative dimension (r-1)
- 1 - m of $P^{\vee}$ i.e. of relative codimension (m+1) in $P^{\vee}$.  (N.B. (r-%
1) is here the relative common dimension of P and $P^{\vee}$ over S), which
we may visualize geometrically \uline{as follows}.  Let us first of all
take $n = r - 1$, we find an isomorphism $\underline{P}(E^{\vee}) \simeq \underline{\mathrm{Grass}}_{r-1}(\underline{E})$ that
allows us to identify the points of $P^{\vee}$ with values in S (let us

\clearpage
% EGA V, English translation by P. Blass, 1987 (cote 24): archive cote 24, site page 74 (PDF page 75).
% typescript. Draft compared: none.
% AI-transcribed by Claude Opus 5.5 (claude-opus-5-5, adaptive thinking, effort high), Message Batches API, batch ega24b (msgbatch_01RGJJeHdR96U16ZWFuUjtDx), 2026-11-13; not yet checked by eye.
\scanpage{24/74}
\mspage{74}

say) as linear subvarieties of codim 1 of P.  (called \uline{again}
\uline{hyperplanes} of P).

This says that $L^0$ consists of hyperplanes which \uline{contain} the
linear subvariety L of P [by which of course we mean] that the
points of L with values in $S^1$ are the hyperplanes in $P_S{}^1$ that
contain $L_S{}^1$).  This follows from the fact that should have occured
at the same time as the fact that a linear subvariety L of $\underline{P}$
determines a locally free quotient G or Q [illegible] of E from
which it results that if G and $G^1$ are two locally free quotients of
E (not necessarily of the same rank) then $\underline{P}(G^1)$ C $\underline{P}(G)$ in so far as
the linear subvarieties of $\underline{P}(E)$, $G^1$ is majorized by G (and the
inclusion $\underline{P}(G^1) \rightarrow \underline{P}(G)$ is nothing else but the deduced morphism G
$\rightarrow G^1$).

Here is a minimum of the [sorites Fr]* which we must have at
our disposal.  The complete list cannot in any case be fixed (only
when) until the sets of other N$^{o1}$s of the present paragraph are
written up.

It seems to me convenient [commode Fr] to introduce also the
functor Grass (E) ($S^1$) = set of quotient modules locally free (of
rank not specified) of $E_S1$ [illegible ask AG]
then \uline{Grass} (E) is representable by n $\geq$ 0 $\mathrm{Grass}_n$(E).  The linear
subvarieties of $\underline{P}$(E) are indeed defined by sections of Grass (E)
over S [NB:  the rank i.e. the relative dimension may vary if S is
not connected.  [slightly illegible confirm with AG].

\margin{* Ask AG or Deligne about the best word.}

\clearpage
% EGA V, English translation by P. Blass, 1987 (cote 24): archive cote 24, site page 75 (PDF page 76).
% typescript. Draft compared: none.
% AI-transcribed by Claude Opus 5.5 (claude-opus-5-5, adaptive thinking, effort high), Message Batches API, batch ega24b (msgbatch_01RGJJeHdR96U16ZWFuUjtDx), 2026-11-13; not yet checked by eye.
\scanpage{24/75}
% NOTE: this side of the sheet has no typed text of its own. The typed text visible on the scan is show-through, in mirror image, from the reverse of the sheet: a French text ("fondation d'un mouvement international ... ``Survivre'' ...") followed by a list of references ("Espaces Vectoriels Topologiques", ``Topologie et algèbre homologique'', ``Géométrie analytique''). It is not transcribed here.
% NOTE: the sheet is torn at the top-left corner, along the right edge and at the lower right.
\hand{\unsure{Idem} cf \unsure{aussi} Groth. pour avoir plusieurs}
\hand{activités \inserted{\uline{importantes}} à la fois, une mathématique,}
\hand{une autre non math. mais valable.}

\mspage{\hand{75}}

\clearpage
% EGA V, English translation by P. Blass, 1987 (cote 24): archive cote 24, site page 76 (PDF page 77).
% typescript. Draft compared: none.
% AI-transcribed by Claude Opus 5.5 (claude-opus-5-5, adaptive thinking, effort high), Message Batches API, batch ega24b (msgbatch_01RGJJeHdR96U16ZWFuUjtDx), 2026-11-13; not yet checked by eye.
\scanpage{24/76}
\mspage{-3-}

fondation d'un mouvement international et interprofessionnel
``Survivre'', et à consacrer aux questions liées à la survie une
part importante de mon énergie. Dans cette optique, la seule va-%
\margin{(}leur de mon apport comme mathématicien est de me permettre aujourd'
hui, grâce à l'estime professionnelle et personnelle acquise parmi
mes collègues, de donner plus de force à mon témoignage et à mon
action en faveur d'une stricte subordination de toutes nos acti-%
vités, y compris nos activités de scientifiques, aux impératifs
\margin{(}de la survie, et à la promotion d'un ordre stable et humain sur
notre planète, sans lequel la survie de notre espèce ne serait
ni possible, ni désirable.

A. GROTHENDIECK.

PRINCIPALES PUBLICATIONS
------------------------

\uline{Espaces Vectoriels Topologiques}

[1] Critères de compacité dans les espaces fonctionnels géné-%
raux, Amer. J. 74 (1952), p. 168-186.

[2] Sur certains espaces de fonctions holomorphes, J. Crelle
192 (1953), p. 35-64 et 77-95.

[3] Espaces Vectoriels Topologiques, Notes polyc., Sa\~o Paulo
(1954), 240 p.

[4] Sur les espaces (F) et (DF), Summa Bras. 3 (1954),
p. 57-123.

[5] Produits tensoriels topologiques et espaces nucléaires,
Mem. AMS, n° 16 (1955), 329 p.

[6] Résumé de la théorie métrique des produits tensoriels
topologiques, Bull. Sa\~o Paulo 8 (1953), p. 1-79.

[7] La théorie de Fredholm, Bull. SMF, 84 (1956), p. 319-384.

\uline{Topologie et algèbre homologique}

[8] Sur quelques points d'algèbre homologique, Tohoku M.J.,
9 (1957), p. 119-221.

[9] Théorèmes de finitude pour la cohomologie des fasceaux,
Bull. SMF, 84 (1956), p. 1-7.

\uline{Géométrie analytique}

[10] Sur la classification des fibrés holomorphes sur la
sphère de Riemann, Amer. J., 79 (1957), p. 121-138.

[11] Techniques de construction en géométrie analytique,
Sem. H. Cartan, 13 (1960/61), exposés 7 à 17.

...

\margin{76}

\end{document}
